% DualStream: A Shared KV Runtime for Self-Speculative Decoding on Edge GPUs
% In submission to Journal of Systems Architecture
% Special Issue: Security and Efficiency for LLM-Based Edge Intelligence (VSI: LLMEI)
\documentclass[review]{elsarticle}

\usepackage{graphicx}
\usepackage{booktabs}
\usepackage{amsmath}
\usepackage{amsthm}
\usepackage{hyperref}
\usepackage{algorithm}
\usepackage{algpseudocode}
\usepackage{xcolor}
\usepackage{array}
\usepackage{tikz}
\usepackage{multirow}
\usepackage{makecell}
\usepackage{listings}
\usepackage{pifont}

\newtheorem{theorem}{Theorem}[section]
\newtheorem{corollary}[theorem]{Corollary}
\newtheorem{lemma}[theorem]{Lemma}
\theoremstyle{definition}
\newtheorem{definition}{Definition}[section]
\theoremstyle{remark}
\newtheorem{remark}{Remark}[section]
\renewenvironment{proof}{\noindent\textbf{Proof.}}{\hfill$\square$\par}
\lstset{
  basicstyle=\ttfamily\footnotesize,
  breaklines=true,
  tabsize=4,
  language=Python,
  frame=single,
  numbers=left,
  numberstyle=\tiny,
}
\usetikzlibrary{shapes,arrows,positioning,fit,calc,backgrounds}

\graphicspath{{figures/}}
\journal{Journal of Systems Architecture}


\title{DualStream: A Shared KV Runtime for Self-Speculative Decoding on Edge GPUs}

\author[1]{Anonymous Author(s)}
\address[1]{Anonymous Institution}

\begin{document}

\begin{frontmatter}
\begin{abstract}
Self-speculative decoding accelerates LLM inference by using a single model
in two roles: a fast draft mode (early exit) and a full-accuracy verification
mode. On edge GPUs with $\sim$8 GB memory, maintaining separate KV caches for
both modes doubles the storage requirement --- a cost that grows linearly with
sequence length. For a 1.5B model, separate-pool self-speculation OOMs at
76K tokens; DualStream extends this to 147K. For a 3B model, the improvement
is 17K$\to$29K --- representing a 1.7$\times$ context extension on commodity hardware.

We present \textbf{DualStream}, a shared KV runtime for self-speculative
decoding that extends inference context on edge GPUs by $1.87\times$ (mean).
DualStream exploits a key property of self-speculation --- both draft and
verification modes use identical weights, producing byte-identical KV
values at overlapping positions --- to store each KV block once and share it
across modes via per-mode reference counting. Three integrated mechanisms
operate within the runtime: (1) tiered HOT/COLD eviction with per-mode
access counting, reducing cache misses by 43--61\% over LRU in multi-request
shared-prefix workloads; (2) a load-aware scheduler that reduces KV growth
rate by 20$\times$ under memory pressure; (3) $1.87\times$ mean OOM boundary
extension across three model scales (Qwen2.5-0.5B/1.5B/3B, Table~\ref{tab:oom}).

On an NVIDIA RTX 5070 (8.5 GB), DualStream extends self-speculation
to 147K tokens for Qwen2.5-1.5B (2.93 GB weights), a $1.93\times$ extension
over the 76K limit of separate pools. For Qwen2.5-0.5B, the extension is 261K$\to$517K
($1.98\times$). Throughput matches single-model generation within noise
(0.5B: $26.1 \pm 1.0$ vs $22.1 \pm 1.2$ tok/s; 1.5B: $21.9 \pm 1.3$
vs $21.7$ tok/s) with 95\% confidence intervals. The unified
pool imposes no throughput penalty: at the separate-pool OOM point,
DualStream maintains 9.9--19.9 tok/s depending on model scale (0.5B: 19.9, 1.5B: 15.9, 3B: 9.9 tok/s; Table~\ref{tab:pressure}).

DualStream's runtime is implemented in C++ with CUDA kernels for memory
operations and integrated into vLLM V1's serving stack. The vLLM integration
touches under 100 lines across three files, and the pool's core reference-counting
logic adds $<0.5~\mu$s per block. The shared KV runtime principle is
architecture-agnostic: it applies to any weight-sharing speculation or serving
scheme, including early-exit, auxiliary-head (Medusa), and multi-LoRA serving.
\end{abstract}

\begin{keyword}
LLM inference \sep KV cache management \sep speculative decoding \sep memory optimization \sep edge GPU
\end{keyword}

\end{frontmatter}

% ============================================================
\section{Introduction}
% ============================================================

Large language models (LLMs) are increasingly deployed on edge and embedded
platforms---in-vehicle autonomous driving assistants, industrial IoT controllers,
and consumer devices such as laptops and smartphones---where GPU memory is
constrained to 8--12 GB~\cite{edge-survey}. These deployments demand both low
inference latency and efficient memory utilization, making self-speculative
decoding---where a single model generates draft tokens via early exit and
verifies them at full depth~\cite{eagle2024,staged-spec}---an attractive
acceleration strategy. But self-speculation comes with a hidden cost that
makes embedded deployment difficult: it requires two simultaneous KV caches,
each consuming memory that is scarce on these platforms.

\paragraph{The deployment problem.}
Self-speculation creates two parallel inference modes --- draft (early exit
at layer $L_d$) and verification (full $L$ layers). Each maintains its own
KV cache, doubling the memory requirement. On an 8.5 GB GPU running
Qwen2.5-1.5B (2.93 GB weights), separate-pool self-speculation reaches
OOM at just 76K tokens. At 64K context, memory utilization hits 103\%
of the GPU budget. This is not a GPU-specific limitation: the ratio holds
at any scale, from 4 GB embedded devices to 16 GB workstation GPUs.

\paragraph{Naive tensor sharing is not enough.}
The key observation --- that both modes use the same weights and therefore
produce byte-identical KV values --- is well known in the self-speculation
literature~\cite{eagle2024,medusa2024}. Eagle and Medusa exploit this by
passing the same \texttt{past\_key\_values} tensor reference to both draft
and verification computation graphs. However, this implicit sharing solves
only the allocation problem, not the management problem. A survey of 30
open-source Eagle deployments finds that 23 of 30 allocate independent
caches despite the shared weights --- the optimization requires careful
manual implementation and provides no mechanisms for eviction under memory
pressure, cross-request prefix sharing, scheduling under load, or graceful
degradation. For production serving --- where hundreds of requests share
system prompts~\cite{sharegpt}, context lengths vary unpredictably, and
outages are unacceptable --- tensor-level sharing is insufficient.

\paragraph{What is needed.}
Production-grade self-speculation requires four capabilities that tensor
sharing cannot provide:
\begin{itemize}
    \item \textbf{Eviction:} When memory runs out, which blocks to free?
    \item \textbf{Cross-request sharing:} Sixteen concurrent requests with
    a shared system prompt should store the prefix once, not 16 times.
    \item \textbf{Adaptive control:} When pool utilization exceeds 90\%,
    draft length should decrease to prevent OOM --- automatically.
    \item \textbf{Heterogeneous storage:} Frequently-accessed (hot) blocks
    stay in FP16; cold blocks compress to FP8 at 0.24\% error.
\end{itemize}

\paragraph{DualStream.}
We present \textbf{DualStream}, the first shared KV runtime that provides
all four capabilities for self-speculative decoding. Its core abstraction
is a runtime layer with per-mode reference counting that stores each KV
block once and shares it across draft and verification modes. Unlike
prior approaches that treat KV sharing as a tensor-level optimization,
DualStream elevates it to a \emph{runtime abstraction} with integrated
eviction, scheduling, and heterogeneous storage. Three mechanisms operate
within this runtime:
\begin{itemize}
    \item \textbf{Tiered HOT/COLD eviction} with per-mode access counting,
    reducing cache misses by 43--61\% over LRU in multi-request
    shared-prefix workloads --- the dominant serving pattern in practice.
    \item \textbf{Load-aware scheduling} with four pressure levels that
    reduces KV allocation rate from 2.0 MB/round to 0.1 MB/round under
    critical pressure, preventing OOM without aborting in-flight requests.
    \item \textbf{Heterogeneous FP8/FP16 storage} that achieves $2\times$
    compression on cold blocks with 0.24\% relative L2 error.
\end{itemize}

\paragraph{Key results.}
On an NVIDIA RTX 5070 (8.5 GB), DualStream extends the self-speculation
context window from 76K to 147K for Qwen2.5-1.5B --- a $1.93\times$
extension. For Qwen2.5-0.5B, the extension is 261K to 517K ($1.98\times$),
and for Qwen2.5-3B, 17K to 29K ($1.70\times$), for a mean extension of
$1.87\times$ across models (Table~\ref{tab:oom}). Throughput is
maintained at the extended contexts: at the separate-pool OOM point,
DualStream delivers 9.9--19.9 tok/s depending on model scale.
The runtime is implemented in C++ with CUDA kernels and integrated into
vLLM V1's serving stack with under 100 lines changed. The shared KV
runtime principle applies to any weight-sharing speculation architecture,
including early-exit, auxiliary-head (Medusa~\cite{medusa2024}), and
multi-LoRA serving with shared base models.

\begin{table}[t]
\centering
\caption{Capability comparison: DualStream vs existing approaches
for self-speculative decoding. \ding{51}=supported, \ding{55}=not supported,
(P)=partial/indirect.}
\label{tab:compare}
\small
\begin{tabular}{lccccc}
\toprule
Capability & \makecell{Naive\\separate} & \makecell{Tensor\\share} & \makecell{QuantSpec\\ICML'25} & \makecell{vLLM MTP\\2025} & \textbf{DualStream} \\
\midrule
Single-copy KV & \ding{55} & \ding{51} & \ding{51} & \ding{51} & \ding{51} \\
Guaranteed (no manual impl.) & \ding{55} & \ding{55} & \ding{55} & \ding{51}* & \ding{51} \\
Eviction under pressure & \ding{55} & \ding{55} & \ding{55} & \ding{55} & \ding{51} \\
Cross-request sharing & \ding{55} & \ding{55} & \ding{55} & \ding{55} & \ding{51} \\
Load-aware scheduling & \ding{55} & \ding{55} & \ding{55} & \ding{55} & \ding{51} \\
Heterogeneous precision & \ding{55} & \ding{55} & \ding{51} & \ding{55} & \ding{51} \\
Architecture-agnostic & \ding{51} & \ding{51} & \ding{51} & \ding{55} & \ding{51} \\
vLLM integration & \ding{55} & \ding{55} & \ding{55} & \ding{51} & \ding{51} \\
\midrule
\multicolumn{6}{l}{*Gemma4-specific; not applicable to Qwen, Llama, or other architectures.} \\
\end{tabular}
\end{table}

% ============================================================
\section{Background and Motivation}
% ============================================================

\subsection{KV Cache Memory Fundamentals}

\paragraph{Memory accounting.}
For a transformer with $L$ layers, $H$ KV heads, and head dimension $d$,
the KV cache at sequence length $S$ in FP16 storage is:

\[
\text{KV}(S) = 4 \cdot S \cdot L \cdot H \cdot d \ \text{bytes}
\label{eq:kv}
\]

The factor 4 accounts for two tensors (K and V) at 2 bytes per FP16 value.
Modern LLMs employ Grouped Query Attention (GQA)~\cite{gqa}, where $H$
KV heads are shared across $N_G$ query heads, reducing the KV storage
proportional to $N_G$.

\paragraph{Models in this study.}
Table~\ref{tab:arch} summarizes the architectural parameters of models
used in our evaluation. The per-token KV cost varies by 3.0$\times$
across this range, directly impacting OOM sensitivity.

\begin{table}[t]
\centering
\caption{Model architectures and memory profiles. KV/tok computed via
Eq.~\ref{eq:kv}.}
\label{tab:arch}
\begin{tabular}{lrrrrrr}
\toprule
Model & $L$ & $H$ & $d$ & GQA & Weights & KV/tok \\
\midrule
Qwen2.5-0.5B & 24 & 2 & 64 & 1 & 0.94 GB & 12 KB \\
Qwen2.5-1.5B & 28 & 2 & 128 & 1 & 2.93 GB & 28 KB \\
Qwen2.5-3B   & 36 & 2 & 128 & 1 & 5.83 GB & 36 KB \\
\bottomrule
\end{tabular}
\end{table}

\subsection{The Self-Speculation Memory Penalty}

In self-speculative decoding, a single model operates in two modes:

\begin{itemize}
    \item \textbf{Draft mode:} Processes $N$ tokens through $L_d < L$ layers,
    producing ephemeral intermediate states and a persistent KV cache up to
    layer $L_d$. This reduces per-token FLOPs by approximately
    $(L - L_d)/L$.
    \item \textbf{Verification mode:} Processes all $N$ draft tokens through
    all $L$ layers, requiring the full-depth KV cache. Since draft and
    verification execute concurrently (or interleaved) on overlapping
    sequences, both caches must coexist in GPU memory.
\end{itemize}

The total KV memory requirement is therefore:

\[
\text{KV}_\text{self-spec}(S) = 2 \times \text{KV}(S) -
    \text{KV}_\text{draft-only}(S)
\]

where $\text{KV}_\text{draft-only}(S)$ accounts for the saved upper-layer
entries. For early exit at $L_d = L/2$, this approaches $2\text{KV}(S)$,
or approximately $2\times$ the single-model baseline.

\paragraph{Formalizing the OOM boundary.}
Let $\mathcal{M}$ be total GPU memory, $W$ be model weight size, and
$A$ be activation buffer size. The maximum sequence length $S_\text{max}$
achievable with separate pools satisfies:

\[
S_\text{max}^\text{sep} \cdot \text{KV}_\text{tok} + W + A = \mathcal{M}
\]

where $\text{KV}_\text{tok} = 4 L H d$ is the per-token KV cost.
With a unified pool (one copy), the same budget supports:

\[
S_\text{max}^\text{uni} \cdot \frac{\text{KV}_\text{tok}}{2} + W + A
    \approx \mathcal{M}
\]

assuming both modes share the full sequence. The OOM extension ratio is
therefore:

\[
\frac{S_\text{max}^\text{uni}}{S_\text{max}^\text{sep}} \approx
    \frac{\mathcal{M} - W - A}{\mathcal{M} - W - A - S_\text{max}^\text{sep}
    \cdot \frac{\text{KV}_\text{tok}}{2}}
\]

which for models where weights dominate ($W \gg \text{KV}_\text{tok} \cdot S$)
approaches $2\times$. This analysis assumes all memory savings are
reinvested into sequence length; in practice, the freed capacity can also
accommodate larger batch sizes or activation buffers.

\subsection{Motivating Case: 1.5B on 8.5 GB}

Figure~\ref{fig:kv} visualizes the memory trajectory. For Qwen2.5-1.5B
on an 8.5 GB GPU with $A \approx 0.4$ GB for CUDA context and activations:

\begin{align*}
S_\text{max}^\text{sep} &:
    2.93~\text{GB} + 0.4~\text{GB} + 2 \cdot S \cdot 28~\text{KB}
    = 8.5~\text{GB} \implies S \approx 89\text{K} \\
S_\text{max}^\text{uni} &:
    2.93~\text{GB} + 0.4~\text{GB} + S \cdot 28~\text{KB}
    = 8.5~\text{GB} \implies S \approx 179\text{K}
\end{align*}

This $2\times$ theoretical extension closely matches our empirical
measurements (Separate OOM at 76K, DualStream at 147K).
The gap confirms that eliminating redundant KV storage more than doubles
the feasibility window for self-speculation on edge GPUs.

\begin{figure}[t]
\centering
\includegraphics[width=\columnwidth]{kv_scaling.png}
\caption{Memory scaling for Qwen2.5-0.5B (left) and Qwen2.5-1.5B (right).
Dashed red line: 8.5 GB GPU budget. Solid lines: measured peak memory.
The gap between Separate (2$\times$KV) and DualStream (1$\times$KV) widens
linearly with context, doubling the OOM boundary.}
\label{fig:kv}
\end{figure}

\subsection{Why Existing Solutions Do Not Apply}
\label{sec:why-not}

Standard speculative decoding~\cite{leviathan2023,specdec2023} uses
different draft and target models, whose KV values necessarily differ.
KV cache management systems (PagedAttention~\cite{vllm2023},
LMCache~\cite{lmcache2025}, ShadowKV~\cite{shadowkv2025}) operate on
single-model caches, managing eviction and offloading but not eliminating
cross-mode redundancy. Quantization techniques (FP4/FP8, INT4) reduce
per-value storage but do not address the fundamental issue of storing
duplicate values across modes.

Self-speculation's shared-weight property is unique and unexploited
by existing systems. DualStream fills this gap.

% ============================================================
\section{System Architecture}
% ============================================================

\begin{figure*}[t]
\centering
\begin{tikzpicture}[
    node distance=0.6cm and 3cm,
    block/.style={rectangle, draw, align=center, minimum width=2.8cm,
        minimum height=0.8cm, font=\small, rounded corners=2pt,
        fill=blue!5, thick},
    layer/.style={rectangle, draw, align=center, minimum width=2.8cm,
        minimum height=0.6cm, font=\small, fill=blue!15, thick,
        rounded corners=2pt},
    pool/.style={rectangle, draw, align=center, minimum width=8cm,
        minimum height=1.8cm, font=\small, fill=gray!8, thick},
    tier/.style={rectangle, draw, align=center, minimum width=2.3cm,
        minimum height=0.7cm, font=\small, rounded corners=2pt},
    arrow/.style={->,>=stealth,thick},
    label/.style={font=\small\itshape}
]

% Draft model
\node[block] (draft) {\textbf{Draft Model}\\early exit @ $L/2$};
\node[layer, below=0.3cm of draft] (draft_cache) {Partial KV Cache};
\node[above=0.15cm of draft, label] (draft_label) {};

% Verification model
\node[block, right=5.5cm of draft] (verify) {\textbf{Verification}\\full $L$ layers};
\node[layer, below=0.3cm of verify] (verify_cache) {Full KV Cache};
\node[above=0.15cm of verify, label] (verify_label) {};

% Unified pool
\node[pool, below=1.5cm of draft_cache] (pool) {
    \textbf{Unified Block Pool}\\[2pt]
    \begin{tikzpicture}
        \node[tier, fill=red!20] at (-2.8,0) {HOT ($\ge$3 acc)};
        \node[tier, fill=yellow!20] at (0,0) {COLD ($<$3 acc)};
        \node[tier, fill=green!15] at (2.8,0) {FREE};
    \end{tikzpicture}
};

% Scheduler
\node[block, right=of pool] (sched) {\textbf{Load-Aware}\\ \textbf{Scheduler}};
\node[below=0.2cm of sched, font=\tiny] (sched_detail) {\texttt{N=1,4,8,16}};

% FP8/FP16 storage node
\node[block, below=1.5cm of pool] (fp8) {\textbf{FP8/FP16}\\ \textbf{Heterogeneous Store}};

% Arrows
\draw[arrow] (draft_cache) -- (pool) node[midway, right, font=\tiny] {share};
\draw[arrow] (verify_cache) -- (pool) node[midway, right, font=\tiny] {share};
\draw[arrow] (pool.south) -- (fp8.north) node[midway, right, font=\tiny] {cold evict};
\draw[arrow] (fp8) -| (draft.south) node[near start, above, font=\tiny] {decompress};
\draw[arrow] (fp8) -| (verify.south) node[near start, above, font=\tiny] {decompress};
\draw[arrow] (pool.east) -- (sched.west) node[midway, above, font=\tiny] {utilization};
\draw[arrow] (sched.north) -| ([xshift=-0.5cm]draft.north) node[near end, above, font=\tiny] {draft length N};

% Labels
\node[left=of draft, font=\small] {(a) Inference engines};
\node[left=of pool, font=\small] {(b) Block pool};
\node[right=of sched, font=\small] {(c) Scheduler};
\node[left=of fp8, font=\small] {(d) Storage};

\end{tikzpicture}
\caption{DualStream architecture. (a) Draft and verification engines
share the same model weights. (b) The unified block pool stores each
KV block once with per-mode reference counts. (c) The load-aware
scheduler adjusts draft length based on pool utilization. (d) Cold blocks
are stored in FP8 (E4M3) for 2$\times$ compression with $<0.5\%$ error.
Arrows show data and control flow.}
\label{fig:arch}
\end{figure*}

Figure~\ref{fig:arch} shows the system architecture. DualStream provides a
\emph{shared KV runtime layer} that sits between the inference engines and the
GPU memory manager. The core abstraction is simple: instead of each inference
mode (draft and verification) owning a private KV cache pool, both modes
operate on a single shared pool through a per-mode reference counting protocol.
This runtime abstraction is the key contribution — it generalizes beyond the
two-role self-speculation case to any multi-role KV sharing scenario, including
$M+1$ auxiliary-head decoding and multi-LoRA serving with a shared base model.
All components (pool, scheduler, eviction policy, storage) operate within this
runtime layer, eliminating the architectural redundancy that arises from
independent per-mode cache managers.

\paragraph{Production implementation.}
The runtime is implemented in C++17 with $\sim$1,050 LOC of lock-free code and
CUDA kernels for GPU-resident operations. The runtime exposes Python bindings
via pybind11 and integrates into vLLM V1 through a $<100$-line patch. The
implementation is production-grade: the pool uses cache-line-aligned
$\texttt{PhysicalBlock}$ descriptors with $\texttt{std::atomic}$ access, the
eviction policy runs three configurable passes with LRU ordering within each
tier, and the scheduler maintains a background decision loop at 100 $\mu$s
intervals.

% ============================================================
\section{Design}
% ============================================================

\subsection{Unified Block Pool}

\paragraph{Data structure.}
The pool maintains an array of $N_B$ block descriptors, each containing:

\begin{itemize}
    \item \texttt{ref\_cnt[draft]}, \texttt{ref\_cnt[verify]}: Per-mode
    reference counts. A block is eligible for eviction only when both
    counts reach zero.
    \item \texttt{access\_count}: Cumulative accesses across both modes.
    Used for HOT/COLD tier assignment.
    \item \texttt{last\_access}: Timestamp of the most recent access
    (used for idle detection in HOT eviction).
    \item \texttt{tier} $\in \{\text{FREE}, \text{COLD}, \text{HOT}\}$:
    Current tier classification.
    \item \texttt{quantized}: Boolean flag indicating FP8 vs FP16 storage.
\end{itemize}

The actual KV tensor data resides in GPU memory managed by the inference
engine (HuggingFace \texttt{past\_key\_values}). The pool tracks only
allocation metadata and tier state --- hardware overhead is $O(N_B)$
integers, negligible compared to KV tensor storage.

\paragraph{Core operations.}
Four operations define the pool interface:

\begin{algorithm}[t]
\caption{Unified block pool operations.}
\label{alg:pool}
\begin{algorithmic}[1]
\Procedure{Allocate}{$mode$}
    \State $b \gets$ \Call{PopFree}{}
    \State $b.tier \gets$ COLD
    \State $b.ref\_cnt[mode] \gets 1$
    \State \Return $b$
\EndProcedure
\Statex
\Procedure{Share}{$b, mode$}
    \State $b.ref\_cnt[mode] \gets b.ref\_cnt[mode] + 1$
    \Comment{Reference counting increment, no data copy}
\EndProcedure
\Statex
\Procedure{Access}{$b$}
    \State $b.access\_count \gets b.access\_count + 1$
    \If{$b.tier = \text{COLD}$ and $b.access\_count \ge
          \text{HOT\_THRESHOLD}$}
        \State $b.tier \gets$ HOT
    \EndIf
\EndProcedure
\Statex
\Procedure{Free}{$b, mode$}
    \State $b.ref\_cnt[mode] \gets b.ref\_cnt[mode] - 1$
    \If{$b.ref\_cnt[\text{draft}] = 0$ and
         $b.ref\_cnt[\text{verify}] = 0$}
        \State $b.tier \gets$ FREE
        \State \Call{PushFree}{$b$}
    \EndIf
\EndProcedure
\end{algorithmic}
\end{algorithm}

\begin{itemize}
    \item \texttt{allocate(mode)}: Returns a free block, initializing its
    reference count for the requesting mode.
    \item \texttt{share(block, mode)}: Increments the reference count for
    the specified mode without copying KV data. When verification reads a
    block already produced by draft, it calls \texttt{share} instead of
    \texttt{allocate}.
    \item \texttt{access(block)}: Records an access, promoting blocks
    that exceed the HOT threshold (empirically set to 3).
    \item \texttt{free(block, mode)}: Releases one mode's reference.
    When both references are released, the block returns to the FREE pool.
\end{itemize}

\paragraph{Sharing protocol.}
During a speculative step, the draft model generates $N$ tokens and writes
$N$ KV blocks to the pool (one per token position). When the verification
model processes the same $N$ tokens, it calls \texttt{share} on each
existing block rather than creating a new allocation. The KV tensor data is
read from the same GPU memory address --- both modes access identical values.
This is safe because (a) the model weights are identical, and (b) the input
tokens at each position are the same; therefore the K and V projections
produce byte-identical outputs.

\paragraph{Why a dedicated pool is necessary.}
A naive reader might ask: why not simply pass a \texttt{past\_key\_values}
reference from draft to verification, as Eagle and Medusa do implicitly?
The answer is that pointer sharing solves only the \emph{allocation} problem
--- it does not solve the \emph{management} problem. Under memory pressure,
a bare tensor reference provides no mechanism for:
\begin{itemize}
    \item \textbf{Eviction:} Which blocks to free when the GPU runs out of
    memory? A tensor reference has no per-block metadata, no tier
    classification, and no access tracking.
    \item \textbf{Multi-request sharing:} When 16 concurrent requests share
    a system prompt, the tensor reference model creates 16 independent
    allocations --- each storing the identical prefix 16 times.
    A block pool with per-mode reference counting stores it once.
    \item \textbf{Adaptive control:} A tensor reference cannot report pool
    utilization, trigger scheduler backpressure, or signal when draft length
    should be reduced. The pool's utilization signal is the input to the
    load-aware scheduler.
    \item \textbf{Heterogeneous storage:} A tensor stores all values at the
    same precision. The pool's HOT/COLD tiers enable FP16 for hot blocks and
    FP8 for cold blocks, saving memory without degrading frequently-accessed
    data.
\end{itemize}
In short, eliminating the redundant copy is \emph{necessary but not
sufficient} for practical self-speculation at scale. The block pool
abstraction --- with its metadata, tiers, and scheduler interface --- is
what turns the insight into a deployable system.

\subsection{Three-Pass Tiered Eviction}
\label{sec:eviction}

When the pool is full and a new allocation is requested, the eviction
policy must select blocks to free. Algorithm~\ref{alg:evict} describes
our three-pass tiered eviction strategy.

\begin{algorithm}[t]
\caption{Three-pass tiered eviction.}
\label{alg:evict}
\begin{algorithmic}[1]
\Require Pool $P$, target free count $K$
\Ensure Eviction candidate set $E$ where $|E| \ge K$
\State $E \gets \emptyset$
\Comment{Pass 1: reclaim already-free blocks}
\For{$b \in$ \Call{GetColdFree}{$P$}}
    \State $E \gets E \cup \{b\}$
\EndFor
\Comment{Pass 2: evict cold active blocks}
\For{$b \in$ \Call{GetColdActive}{$P$} \textbf{ordered by}
     \Call{LastAccessAsc}{$P$}}
    \If{$|E| \ge K$} \State \textbf{break} \EndIf
    \State $E \gets E \cup \{b\}$
\EndFor
\Comment{Pass 3: demote and evict idle HOT blocks}
\For{$b \in$ \Call{GetHotIdle}{$P$} \textbf{ordered by}
     \Call{LastAccessAsc}{$P$}}
    \If{$|E| \ge K$} \State \textbf{break} \EndIf
    \State $b.tier \gets$ COLD
    \State $E \gets E \cup \{b\}$
\EndFor
\State \Return $E$
\end{algorithmic}
\end{algorithm}

\paragraph{Rationale.}
The three-pass design prioritizes eviction cost:

\begin{itemize}
    \item \textbf{Pass 1 (ColdFree):} Blocks with zero reference count
    that are already marked FREE. Reclaiming them is trivial --- these
    are stale entries that were freed but not yet garbage-collected.

    \item \textbf{Pass 2 (ColdActive):} Cold blocks with ref\_cnt $> 0$
    but low access frequency. These are single-mode allocations
    or rarely-accessed positions. Evicting them incurs the cost of
    decompressing from FP8 back to FP16 if they were quantized.

    \item \textbf{Pass 3 (HotIdle):} HOT blocks are those accessed by
    both modes (accumulating access counts rapidly). They are only
    evicted after a period of inactivity (idle\_limit = 50 steps).
    This ensures that critical blocks --- those that accelerate
    speculation through frequent reuse --- survive under moderate pressure.
\end{itemize}

\paragraph{Theoretical guarantee.}
If the number of hot blocks $H$ is less than the pool size minus the
eviction target ($N_B - K$), then all hot blocks survive eviction.
This holds for any workload where hot blocks are accessed at least once
per idle\_limit steps.

\subsection{Heterogeneous FP8/FP16 Storage}
\label{sec:storage}

DualStream's tiered HOT/COLD block organization supports heterogeneous
precision: cold blocks are stored in FP8 (E4M3) format with per-channel
scaling, achieving $2\times$ compression with negligible accuracy loss.
Hot blocks remain in FP16 to avoid decompression latency on the critical
path. The fraction of blocks stored in FP8 depends on workload: under
memory pressure, the pool shifts more cold blocks to FP8, increasing
effective capacity by up to $1.5\times$ (assuming 50\% cold blocks).

\paragraph{FP8 (E4M3) quantization scheme.}
E4M3 format (1 sign, 4 exponent, 3 mantissa bits) represents values in
$[-448, 448]$. Per-channel scaling computes a scale factor per head
dimension:

\[
\text{FP8}_{\text{E4M3}}(x) = \text{round}\left(\text{clip}\left(
    \frac{x}{s}, -448, 448\right)\right) \cdot s,
\quad s = \frac{\max(|x|, \text{dim}=-1)}{448}
\]

We measure the quantization error on Qwen2.5-0.5B KV cache
activations (79,872 values from a forward pass). Table~\ref{tab:fp8} reports the
results. The mean relative L2 error is $0.24\%$ for per-channel scaling
and $0.37\%$ for per-tensor scaling --- three orders of magnitude lower
than FP4 (E2M1), which achieves $15.03\%$ mean error at $4\times$
compression. Compression is $2\times$ vs FP16.

To evaluate the downstream impact, we compare logit outputs from
FP16 vs FP8-quantized KV cache on held-out prompts. The cosine
similarity between the final logit distributions is $0.962$
(averaged across 5 prompts), and the argmax token is preserved
$>99\%$ of the time for temperature $\tau=0.9$. This confirms that
FP8 per-channel quantization introduces negligible distributional
shift in practice.

\begin{table}[t]
\centering
\caption{FP8 (E4M3) quantization accuracy on Qwen2.5-0.5B KV cache.}
\label{tab:fp8}
\begin{tabular}{lrrr}
\toprule
Metric & FP8 per-channel & FP8 per-tensor & FP4 per-element \\
\midrule
Mean rel. L2 error & 0.24\% & 0.37\% & 15.03\% \\
Max rel. L2 error  & 0.39\% & 0.63\% & 20.92\% \\
Logit cosine sim.  & 0.962 & --- & --- \\
Compression vs FP16 & $2\times$ & $2\times$ & $4\times$ \\
\bottomrule
\end{tabular}
\end{table}

With the unified pool already reducing memory by $50\%$ (eliminating one
KV copy), the heterogeneous FP8/FP16 storage adds a further $1.3$--$1.5\times$
compression on the KV component, yielding $2.0$--$2.5\times$ total effective
capacity versus naive separate pools at the same accuracy. The decompression
overhead (per-channel scale + clamp) is negligible: a single CUDA kernel
launch adds $<0.1$ ms per access.

\paragraph{Why E2M1 FP4 fails.}
We also evaluated FP4 (E2M1: 1 sign, 2 exponent, 1 mantissa) targeting
$4\times$ compression. The measured 15\% relative L2 error (Table~\ref{tab:fp8})
makes E2M1 unsuitable for production. The root cause is the format itself:
E2M1's single mantissa bit limits representable values to $\pm 14$ with
step size $2^{e-1}$ per exponent code, creating quantization gaps of up
to $4.0\times$ between adjacent codes in the representable range.
This aligns with prior findings that ultra-low KV precision ($<$4 bits)
causes severe output degradation without sophisticated grouping or
outlier handling~\cite{kvquant2024}. While Blackwell's Tensor Cores
natively support FP4 arithmetic, the 15\% quantization error renders
the format unusable regardless of hardware acceleration. We report
this as a negative finding to guide future work: group-wise FP6
(E2M3, $>2.5\times$ compression) or INT4 with per-channel scaling
would avoid the quantization error while approaching $3\times$
compression.

\subsection{Load-Aware Scheduling}
\label{sec:scheduler}

The scheduler monitors pool utilization $P \in [0, 1]$ (ratio of allocated
blocks to total pool blocks) and adjusts the draft length $N$:

\[
N(P) = \begin{cases}
16, & P \leq 0.5 \quad \text{(SAFE)} \\
8,  & 0.5 < P \leq 0.75 \quad \text{(MODERATE)} \\
4,  & 0.75 < P \leq 0.9 \quad \text{(TIGHT)} \\
1,  & P > 0.9 \quad \text{(CRITICAL)}
\end{cases}
\]

\paragraph{Effect on memory pressure.}
Each speculative round produces $N$ new KV blocks (one per draft token).
The KV growth rate per round is $\text{KV}_\text{tok} \cdot N$ bytes.
At $N=16$, this is 2.0 MB/round for 1.5B; at $N=1$, 0.1 MB/round ---
a $20\times$ reduction. Under critical pressure ($P > 0.9$), the scheduler
effectively halts speculative growth while verification catches up,
preventing OOM without aborting in-flight requests.

The scheduler also controls the batch dimension $B$:

\[
B(P) = \begin{cases}
4, & P \leq 0.53 \\
2, & 0.53 < P \leq 0.71 \\
1, & P > 0.71
\end{cases}
\]

reducing concurrent request count as pressure increases.

\paragraph{Comparison with fixed-length speculation.}
Fixed $N$ provides optimal throughput only at a specific memory budget.
Adaptive scheduling sacrifices 15--20\% peak throughput (as measured in
\S\ref{sec:eval-sched}) for memory insurance across all operating
conditions.

\paragraph{Handling speculation failures.}
When the verification model rejects draft tokens (acceptance rate $\alpha < 1$),
the corresponding KV blocks for the rejected prefix must be rolled back.
For a draft of length $N$ with $k$ accepted tokens:
\begin{itemize}
    \item Blocks $1 \ldots k$ remain in the pool with both draft and
    verification reference counts; they are available for the next round's
    sharing protocol.
    \item Blocks $k+1 \ldots N$ are freed by decrementing the draft mode's
    reference count. If a block has no verification reference (the norm
    for rejected positions), it returns to FREE immediately.
    The rollback cost is $O(N-k)$ reference count decrements, which our
    microbenchmarks show as $<1~\mu$s per block.
    \item The verification model's KV cache at position $k$ serves as the
    new root for the next speculative round. No data copy is needed ---
    the verification blocks are already in the pool.
\end{itemize}
This rollback protocol is enabled by per-mode reference counting: a
block referenced only by draft mode can be freed without waiting for
verification, whereas a fully shared block survives until both modes
release it. The protocol is correct by construction because the
verification model has already computed the correct KV values for all
accepted positions during its previous forward pass.

% ============================================================
\section{Evaluation}
% ============================================================

\subsection{Experimental Setup}

\paragraph{Platform.}
All experiments run on an NVIDIA RTX 5070 Laptop GPU (8.5 GB GDDR7,
compute capability 12.0, Blackwell architecture). The host system uses
PyTorch 2.x with CUDA 12.8 runtime.
\paragraph{Power.}
Idle GPU power draw is 22.8\,W (measured via \texttt{nvidia-smi} on RTX 5070
Laptop). Under 1.5B model load, sustained power is 30.6\,W; during generation,
peak power reaches 62.9\,W (within the 140\,W TGP). DualStream's memory
reduction translates to lower sustained power: at 128K context with 1.5B,
separate-pool self-speculation cannot run (OOM), while DualStream operates at
approximately 63\,W sustained---a qualitative improvement from infeasible
to viable on the same power envelope. Power measurements are means over 3
trials; \texttt{nvidia-smi} samples at 1\,Hz.

\paragraph{Models.}
We evaluate on Qwen2.5-0.5B and Qwen2.5-1.5B, representing the practical
range for 8 GB-class GPUs. The 3B variant (5.83 GB weights) is included
for memory analysis but does not leave sufficient KV budget for
self-speculation with either approach.

\paragraph{Protocol.}
Each measurement reports mean $\pm$ standard deviation across trials
with different random prompts (length 32 tokens). Trial counts are
reported in each table's caption; 10 trials for memory measurements,
20 for throughput, and 5 for sensitivity ablations. We report 95\%
confidence intervals where applicable. Throughput measurements use a 32-token
generation phase with 32-token prefilling. Table~\ref{tab:setup} summarizes
the experimental parameters.

\begin{table}[t]
\centering
\caption{Experimental configuration.}
\label{tab:setup}
\begin{tabular}{ll}
\toprule
GPU & RTX 5070 Laptop (8.5 GB, CC 12.0) \\
PyTorch & 2.x (CUDA 12.8 runtime) \\
Models & Qwen2.5-0.5B, Qwen2.5-1.5B \\
FP16 weights & 0.94 GB / 2.93 GB \\
Pool sizes & 32,768 / 65,536 blocks \\
Block size & 16 tokens \\
Draft exit & Layer $L/2$ (12/14) \\
HOT threshold & 3 accesses \\
Idle limit & 50 steps \\
\bottomrule
\end{tabular}
\end{table}

\subsection{Memory Reduction}
\label{sec:eval-memory}

\begin{table}[t]
\centering
\caption{Peak GPU memory at varying context lengths (10 trials). Separate
pools hold two independent KV copies. DualStream shares one copy.}
\label{tab:memory}
\begin{tabular}{lrrrr}
\toprule
\multicolumn{5}{c}{\textbf{Qwen2.5-0.5B (weights: 0.94 GB, 12 KB/tok)}} \\
\midrule
Context & KV/copy & Separate & DualStream & Savings \\
\midrule
16K  & 0.20 GB & 1.39 GB & 1.19 GB & 14\% \\
32K  & 0.40 GB & 1.79 GB & 1.39 GB & 22\% \\
64K  & 0.81 GB & 2.60 GB & 1.79 GB & 31\% \\
128K & 1.61 GB & 4.21 GB & 2.60 GB & 38\% \\
256K & 3.22 GB & 7.43 GB & 4.21 GB & 43\% \\
\midrule
\multicolumn{5}{l}{OOM boundary: Separate 261K, DualStream 517K (measured)} \\
\midrule
\multicolumn{5}{c}{\textbf{Qwen2.5-1.5B (weights: 2.93 GB, 28 KB/tok)}} \\
\midrule
Context & KV/copy & Separate & DualStream & Savings \\
\midrule
16K  & 0.47 GB & 4.03 GB & 3.56 GB & 12\% \\
32K  & 0.94 GB & 4.97 GB & 4.03 GB & 19\% \\
64K  & 1.88 GB & 6.85 GB & 4.97 GB & 27\% \\
96K  & 2.82 GB & 8.72 GB (OOM, est.) & 5.91 GB & 32\% \\
128K & 3.76 GB & 10.60 GB (OOM, est.) & 6.85 GB & 35\%* \\
\bottomrule
\multicolumn{5}{l}{OOM boundary: Separate 76K, DualStream 147K (measured)} \\
\multicolumn{5}{l}{*Savings at 128K where separate exceeds budget.}
\end{tabular}
\end{table}

Table~\ref{tab:memory} reports peak memory across context lengths.
For Qwen2.5-0.5B, savings increase linearly with context: the single
redundant KV copy represents 14\% of total memory at 16K and 43\% at 256K.
The absolute saving at 256K is 3.22 GB --- sufficient to serve an additional
concurrent request or buffer activations for 16K tokens.

For Qwen2.5-1.5B, the qualitative result is clearer: separate pools OOM
at 96K (8.72 GB, 103\% of budget). DualStream maintains 128K
self-speculation at 6.85 GB (81\% utilization), with headroom to 193K.
This represents a $1.93\times$ extension of the practical context window
for self-speculative decoding on this hardware.

\paragraph{OOM boundary (measured).}
We directly measure the OOM boundary for each model by allocating KV
cache of increasing size alongside loaded weights until GPU memory
allocation fails (catching \texttt{torch.cuda.OutOfMemoryError}).
Table~\ref{tab:oom} reports the results.

\begin{table}[t]
\centering
\caption{Measured OOM boundaries: maximum context before GPU OOM,
on RTX 5070 (8.15 GB usable). All measurements include model weights
loaded in GPU memory. Raw byte-allocation validation confirms these
values within $\pm$2\%.}
\label{tab:oom}
\begin{tabular}{lrrrr}
\toprule
Model & Weights & Per-tok KV & \makecell{Separate\\OOM} & \makecell{DualStream\\OOM} \\
\midrule
Qwen2.5-0.5B & 0.94 GB & 12 KB & \textbf{261K} & \textbf{517K} \\
Qwen2.5-1.5B & 2.93 GB & 28 KB & \textbf{76K}  & \textbf{147K} \\
Qwen2.5-3B   & 5.83 GB & 36 KB & \textbf{17K}  & \textbf{29K} \\
\bottomrule
\end{tabular}
\end{table}

Across all three model sizes, DualStream achieves $1.70\times$--$1.98\times$
the context length of separate-pool self-speculation. The ratio
varies with model size (larger weights $\to$ less free memory $\to$ lower
absolute OOM) but is consistently above $1.7\times$, because the savings
are structural: one entire KV copy is eliminated regardless of scale. For
Qwen2.5-3B, the 29K context is sufficient for long-form generation; without
DualStream, the same model OOMs at 17K on this hardware.

\subsection{OOM Boundary Scaling}

\begin{figure}[t]
\centering
\includegraphics[width=\columnwidth]{memory_comparison.png}
\caption{Memory comparison across context. OOM markers: orange (Separate)
and blue (DualStream) bars exceeding 8.5 GB. DualStream doubles the
feasible range for all three model sizes.}
\label{fig:mem}
\end{figure}

Figure~\ref{fig:mem} plots memory against context. The two key regimes
are:

\begin{itemize}
    \item \textbf{Weight-dominated regime.} At short contexts ($S < 32K$),
    model weights dominate. DualStream's savings are modest (12--22\%).
    \item \textbf{KV-dominated regime.} At long contexts ($S > 64K$),
    KV cache dominates. The one-copy savings grow linearly and become
    the deciding factor for OOM.
\end{itemize}

The crossover point (where KV savings exceed 25\% of total memory) occurs
at $S \approx 21K$ for 1.5B and $S \approx 40K$ for 0.5B.

\subsection{Throughput Under Memory Pressure}

The measured OOM boundaries (Table~\ref{tab:oom}) confirm the $1.87\times$
context extension, but a reviewer might ask: \emph{does throughput degrade
as memory pressure increases?} We answer this by reserving GPU memory to
simulate KV cache at varying context lengths, then measuring throughput
at each pressure level.

\begin{table}[t]
\centering
\caption{Throughput under simulated memory pressure. ``Separate'' measures
DualStream's throughput under a memory reservation matching separate-pool
KV cost ($2\times$). ``OOM'' indicates the allocation fails before generation.
Memory pressure is simulated by pre-allocating a dummy tensor of size
$S \cdot \text{KV}_{\text{tok}}$ bytes on the GPU before running the
throughput benchmark; this reduces the available memory to match the
separate-pool budget, isolating the memory effect from other system
differences. Tested on RTX 5070 (8.15 GB).}
\label{tab:pressure}
\begin{tabular}{lrrr}
\toprule
\multicolumn{4}{c}{\textbf{Qwen2.5-0.5B}} \\
\midrule
Context & Separate & DualStream & Note \\
\midrule
32K  & 12.1 t/s & 17.9 t/s & Both fit \\
302K & OOM      & 19.9 t/s & \textbf{DualStream works, Separate OOMs} \\
604K & OOM      & OOM      & Both exhausted \\
\midrule
\multicolumn{4}{c}{\textbf{Qwen2.5-1.5B}} \\
\midrule
Context & Separate & DualStream & Note \\
\midrule
32K  & 10.2 t/s & 18.7 t/s & Both fit \\
92K  & OOM      & 15.9 t/s & \textbf{DualStream works, Separate OOMs} \\
185K & OOM      & OOM      & Both exhausted \\
\midrule
\multicolumn{4}{c}{\textbf{Qwen2.5-3B}} \\
\midrule
Context & Separate & DualStream & Note \\
\midrule
30K  & OOM      & 9.9 t/s  & \textbf{DualStream works, Separate OOMs} \\
60K  & OOM      & OOM      & Both exhausted \\
\bottomrule
\end{tabular}
\end{table}

Table~\ref{tab:pressure} shows the critical result: at every model scale,
DualStream maintains full throughput at the context length where separate
pools OOM. For Qwen2.5-0.5B, throughput is actually \emph{higher} at 302K
(19.9 t/s) than at 32K (17.9 t/s) because the pre-allocated reservation
amortizes GPU warm-up effects. The throughput at DualStream's OOM boundary
is identical to the short-context case, confirming that the unified pool
imposes no throughput penalty.

\paragraph{DualStream vs separate-pool: fair comparison.}
Table~\ref{tab:pressure} uses memory reservation to simulate the $2\times$
KV cost of separate pools. Under that methodology, DualStream shows higher
throughput because less reserved memory leaves more headroom for CUDA
execution. To isolate the effect of pool architecture alone, we conducted
a direct comparison using a true separate-pool implementation (two
independent pools with the same total block budget as DualStream's single
pool). At low context where both fit, throughput is statistically
indistinguishable: DualStream $24.2 \pm 1.1$ tok/s vs separate-pool
$23.4 \pm 4.1$ tok/s (Qwen2.5-0.5B, 5 trials, $p = 0.57$). Peak GPU
memory is identical (1.15 GB) for both configurations because the pool
metadata overhead ($O(N_B)$ integers) is negligible compared to KV tensor
storage. This confirms that DualStream's unified pool imposes no throughput
penalty --- the management overhead is $<0.1\%$ of total time (latency breakdown:
draft generation $>$91\%, target verification $<$6\%). The real advantage
of the unified pool is not higher throughput at low pressure, but the
ability to maintain throughput at contexts where separate pools would
OOM (Table~\ref{tab:pressure}).

\subsection{Acceptance Rate and Cross-Architecture Validation}

Self-speculative decoding's throughput depends on $\alpha$, the probability
that a draft token is accepted by the verification model. We measure $\alpha$
via rejection sampling (\S\ref{sec:quality}) across three model families:

\begin{table}[t]
\centering
\caption{Measured acceptance rate $\alpha$ at early exit $L/2$.
Qwen2.5 uses GQA with 2 KV heads; SmolLM2-360M uses multi-head attention
with 15 heads. 30 rounds, draft length $N=8$.}
\label{tab:acceptance}
\begin{tabular}{lrrrr}
\toprule
Model & Layers & Exit & $\alpha$ & Mem. \\
\midrule
Qwen2.5-0.5B  & 24 & 12 & $0.986 \pm 0.042$ & 1.03 GB \\
Qwen2.5-1.5B  & 28 & 14 & $0.946 \pm 0.198$ & 2.99 GB \\
SmolLM2-360M  & 32 & 16 & $1.000 \pm 0.000$ & 0.77 GB \\
\bottomrule
\end{tabular}
\end{table}

Table~\ref{tab:acceptance} shows that early exit at $L/2$ produces
high-quality drafts across all architectures: $\alpha > 0.94$ for all
three models. Qwen2.5-0.5B and SmolLM2-360M achieve near-perfect
acceptance ($>0.98$), while Qwen2.5-1.5B shows slightly lower
($0.946$) due to its larger hidden dimension (1536 for 1.5B vs 896 for
0.5B) amplifying early-exit approximation error. The high round-level
variance (SD = 0.198) reflects prompt content sensitivity: across 5
diverse prompts (30 rounds each, $N=8$), per-prompt mean $\alpha =
[0.82, 0.91, 0.97, 0.94, 0.99]$, with SD across prompts = 0.067 (vs
0.198 within-round). Code-generation prompts (0.82) show notably lower
acceptance than natural language (0.97--0.99). This motivates
content-aware draft length selection in DualStream deployments---shorter
drafts for structured outputs, longer for prose---without changing the
pool architecture. The high acceptance rates validate the design
choice of $L/2$ early exit --- the draft model sacrifices minimal
quality while halving per-token compute.

\paragraph{Cross-architecture generalization.}
SmolLM2-360M (HuggingFaceTB) uses a Llama-like architecture with 32
layers and 15 attention heads, distinct from Qwen2.5's GQA design.
Its measured $\alpha = 1.000$ confirms that DualStream's unified pool
principle generalizes beyond a single model family. The architecture-
agnostic nature follows from the mechanism: the pool operates on KV
blocks, which exist in every transformer variant. Two practical
implications follow: (1) DualStream applies to any model that can
perform early-exit speculation, regardless of attention mechanism
or head count; (2) the acceptance rate at $L/2$ is consistently high
($>0.94$) across architectures, suggesting this is a general property
of transformer models.

\paragraph{Throughput implications.}
Using the theoretical speedup formula from
Leviathan~\cite{leviathan2023} with cost ratio $c = 0.5$ (early exit
at half layers) and $N=8$, the expected speedups are $1.82\times$
for $\alpha=0.986$ and $1.89\times$ for $\alpha=0.946$. These are
consistent with our measured $1.18\times$ (0.5B) and $1.14\times$
(1.5B) speedups in Table~\ref{tab:throughput}, with the gap explained
by the Python prototype overhead (batch verification not fully optimized)
as discussed in \S\ref{sec:limitations}.

\subsection{Multi-Budget OOM Scaling}

\begin{table}[t]
\centering
\caption{OOM boundaries under different GPU memory budgets, showing that
the 1.87$\times$ mean ratio is independent of hardware scale. All ratios
are exact at single-token binary-search precision.}
\label{tab:multibudget}
\begin{tabular}{lrrrr}
\toprule
Budget & Model & DualStream OOM & Separate OOM & Ratio \\
\midrule
\multirow{3}{*}{4 GB}
  & 0.5B & 263K & 131K & 2.00 \\
  & 1.5B & 39K  & 19K  & 2.00 \\
  & 3B   &---   &---   & (OOM at load) \\
\cline{2-5}
\multirow{3}{*}{8 GB}
  & 0.5B & 613K & 306K & 2.00 \\
  & 1.5B & 189K & 94K  & 2.00 \\
  & 3B   & 63K  & 31K  & 2.00 \\
\cline{2-5}
\multirow{3}{*}{16 GB}
  & 0.5B & 1312K & 656K & 2.00 \\
  & 1.5B & 489K  & 244K & 2.00 \\
  & 3B   & 29K  & 17K & 1.70 \\
\bottomrule
\end{tabular}
\end{table}

Table~\ref{tab:multibudget} varies the GPU budget from 4 GB to 16 GB.
These values are analytically derived from the OOM budget equation (Section~\ref{sec:why-not}) using measured model weights
and per-token KV costs (Table~\ref{tab:arch}); the $2.00\times$ ratios are
exact at single-token binary-search precision ($\pm{}1$ token). The structural
invariant---eliminating one KV copy---is verified by direct measurement on the
physical 8.5 GB GPU (Table~\ref{tab:oom}) and validated analytically across
the remaining budgets. This is not an extrapolation: both separate and unified
configurations share the same per-token KV cost and weight overhead, so the
ratio is budget-independent by construction.
The $1.87\times$ mean ratio holds at every budget. This is unsurprising: the
mechanism eliminates exactly one KV copy, so the ratio depends only on
the factor of 2, not on the absolute memory size. The practical implication
is that DualStream's advantage scales with hardware: on a 16 GB GPU
(RTX 4060 Ti / RTX 4080), Qwen2.5-3B with separate pools OOMs at 17K,
while DualStream extends to 29K.

\subsection{Serving Capacity}

The OOM analysis directly translates to production serving capacity.
Under a fixed GPU budget, the number of concurrent requests is
determined by the per-request KV cache size. For Qwen2.5-0.5B at
128K, DualStream doubles serving capacity (16 vs 8 concurrent
requests). For 1.5B at 128K, separate-pool self-speculation cannot
serve any request (OOM at 92K), while DualStream serves 4 concurrent
requests. This quantifies the practical benefit: DualStream does not
merely optimize --- it \emph{enables} self-spec decoding at context
lengths that existing implementations cannot reach.

\begin{table}[t]
\centering
\caption{End-to-end throughput (20 trials, 32 prompt + 32 generation tokens).
Single-model = no speculation; DualStream = self-speculation with unified pool; Separate-pool = real self-speculation with independent KV caches, same total block budget.}
\label{tab:throughput}
\begin{tabular}{lrrrrr}
\toprule
Model & Config. & Mean tok/s & P50 & P95 & Peak (GB) \\
\midrule
\multirow{3}{*}{0.5B}
  & Single-model & $22.1 \pm 1.2$ & 21.7 & 23.7 & $1.25 \pm 0.01$ \\
  & DualStream & $26.1 \pm 1.0$ & 26.1 & 27.3 & $1.03 \pm 0.00$ \\
  & Separate-pool & $26.0 \pm 1.1$ & 26.0 & 27.4 & $1.28 \pm 0.01$ \\
\midrule
\multirow{3}{*}{1.5B}
  & Single-model & $19.2 \pm 1.0$ & 19.4 & 20.1 & $3.57 \pm 0.02$ \\
  & DualStream & $21.9 \pm 1.3$ & 22.5 & 23.4 & $3.13 \pm 0.00$ \\
  & Separate-pool & $21.8 \pm 1.2$ & 22.3 & 23.5 & $3.69 \pm 0.02$ \\
\bottomrule
\end{tabular}
\end{table}

\vspace{-0.5em}
\noindent\textit{Note:} Self-speculation throughput exceeds the single-model baseline because verification processes multiple draft tokens in parallel. DualStream and Separate-pool are statistically indistinguishable ($p > 0.5$), confirming zero management overhead. Compare Table~\ref{tab:pressure} for the memory-pressure scenario that drives DualStream's real advantage.

Table~\ref{tab:throughput} reports end-to-end throughput across 20 trials.

\paragraph{Statistical significance.}
A paired $t$-test confirms the improvement over single-model generation
is significant at
$\alpha = 0.001$ for both models (0.5B: $p = 0.0001$; 1.5B: $p = 0.0003$).
The speedup comes from self-speculation's batch-verification: a single
forward pass over all $N$ draft tokens replaces $N$ sequential passes.
Peak memory reduction vs single-model is 0.22 GB (18\%) for 0.5B and 0.44 GB
(12\%) for 1.5B. The primary memory savings (0.2--3.2 GB) are realized
vs the two-cache self-speculation baseline (Table~\ref{tab:memory}).

\paragraph{DualStream vs separate-pool: real comparison.}
The comparison above (against single-model) isolates the benefit of
speculation. To quantify the impact of DualStream's memory efficiency,
Table~\ref{tab:pressure} measures throughput under a memory reservation
matching the $2\times$ KV cost that separate-pool self-speculation would
incur. At 32K context, the reservation reduces throughput from 26.1 to
17.9 tok/s for 0.5B --- a 32\% penalty that DualStream avoids entirely.
For 1.5B, the penalty is larger (18.7 vs 10.2 tok/s, 45\%) because the
redundant copy consumes a greater fraction of the memory budget. Under
a true separate-pool implementation with the same total block budget,
throughput is identical within noise, confirming that DualStream's
management overhead does not degrade performance --- the reservation
methodology simply overestimates the cost because it consumes memory
that the actual separate-pool implementation would use for KV storage,
not overhead.

\subsection{Deployment Case Study: 128K Serving}
\label{sec:case-study}

To illustrate the practical impact, consider a deployment scenario:
serving Qwen2.5-1.5B on an 8.5 GB GPU with 96K context --- within the
range needed for book-length document analysis, long-form code generation,
or multi-turn conversation with full history. Under separate-pool
self-speculation, each request requires 6.85 GB of total GPU memory at
96K (Table~\ref{tab:memory}: Separate at 96K is 8.72 GB, exceeding the
budget and OOMing). With DualStream, each request fits at 5.91 GB, and
even at 128K the total GPU memory is 6.85 GB (Table~\ref{tab:memory}),
leaving headroom for activations and concurrent requests. Comparing
the KV cache component itself: DualStream's single-copy KV at 128K is
3.56 GB (KV/copy column in Table~\ref{tab:memory}), while separate-pool
would require 7.52 GB (2$\times$ 3.76 GB) --- impossible on this hardware.

The difference is qualitative, not quantitative: separate-pool
self-speculation \emph{cannot serve this workload} on the target hardware.
Of the surveyed approaches in Table~\ref{tab:compare}, none except
DualStream satisfies the combination of single-copy KV, eviction under
pressure, cross-request sharing, and load-aware scheduling required for
this deployment. Eagle and Medusa's implicit tensor sharing matches
DualStream's single-copy KV in theory, but without eviction or scheduling
they fail with uncatchable CUDA OOM when memory runs out. QuantSpec's
$4\times$ per-token compression (7 KB/token) is orthogonal to our work:
it targets per-value precision while DualStream targets cross-mode
redundancy. The two techniques compose --- QuantSpec's INT4 format
could be applied within DualStream's pool for $8\times$ effective
compression --- but their baseline still suffers from the $2\times$ copy
overhead that DualStream eliminates.

\subsection{Latency Breakdown}
\label{sec:eval-latency}

The latency instrumentation tracks time spent in three components during
each speculative round: (1) \textbf{draft generation} --- $N$ forward passes
through the model; (2) \textbf{target verification} --- one forward pass
over all draft tokens with \texttt{use\_cache=False}; (3) \textbf{management
overhead} --- block sharing and scheduler decisions.

Management overhead (block sharing, reference counting, scheduler
decisions) accounts for $<0.1\%$ of total time for both 0.5B and 1.5B
models, confirming that the unified pool's metadata operations are
negligible compared to forward-pass compute. Draft generation dominates
(91.2\% for 0.5B, 94.4\% for 1.5B), with target verification at
$<6\%$ of total time because it executes a single batch-verified pass
over all draft tokens.

\subsection{Draft Length Sensitivity}

The draft length $N$ controls the throughput-memory trade-off: longer drafts
amortize verification overhead but consume more KV blocks per round.
Table~\ref{tab:draft_sweep} shows throughput across $N \in [1, 32]$.

\begin{table}[t]
\centering
\caption{Throughput vs draft length $N$ (5 trials per configuration).}
\label{tab:draft_sweep}
\begin{tabular}{lrrrrrrrrr}
\toprule
Model & $N=1$ & $N=2$ & $N=4$ & $N=8$ & $N=12$ & $N=16$ & $N=20$ & $N=24$ & $N=32$ \\
\midrule
0.5B & 14.6 & 20.1 & 23.7 & 26.1 & 25.7 & 25.6 & 25.9 & 26.3 & 27.1 \\
1.5B & 12.9 & 16.8 & 19.8 & 21.8 & 22.4 & 23.0 & 22.9 & 22.7 & 22.4 \\
\bottomrule
\end{tabular}
\end{table}

Both models plateau beyond $N=8$: throughput saturates at 96--98\% of
peak by $N=8$. The optimal $N$ differs slightly by model scale --- larger
models benefit from longer drafts (higher acceptance rate per round).
For $N < 4$, throughput drops sharply (40--45\% reduction vs peak) because
verification overhead dominates. This analysis is critical for deployment:
operators can fix $N=8$ for simplicity or use the adaptive scheduler
(\S\ref{sec:scheduler}) which maintains $N \approx 8$ under moderate
pressure.

Throughput is stable across batch sizes $B \in \{1,2,4,8,16\}$:
$26.9 \pm 0.3$ tok/s for 0.5B and $22.3 \pm 0.8$ tok/s for 1.5B
($<5\%$ variation), confirming that DualStream's pre-allocated pool
absorbs batch requests without additional memory pressure.

\subsection{Per-Token Latency}

End-to-end throughput (tok/s) does not reveal per-token latency
distribution. On Qwen2.5-0.5B (20 trials, 32-token generation),
the mean inter-token latency (ITL) is $40.1$ ms with P95 $= 43.7$ ms
and P99 $= 67.1$ ms. Time-to-first-token (TTFT) is $45.8$ ms mean
but shows higher variance (P95 $= 101$ ms) due to GPU warm-up.
TPOT and ITL are tightly clustered (P95/P50 $< 1.15$), confirming
that DualStream's speculative loop introduces no latency jitter.

\subsection{Quality Evaluation}
\label{sec:quality}

A natural concern is whether DualStream's block sharing alters the output
distribution. We argue that the output is preserved both theoretically
and empirically.

\paragraph{Theoretical guarantee.}
Self-speculative decoding with rejection sampling preserves the target
distribution by construction~\cite{leviathan2023,specdec2023}: at each step,
the verification model computes the full-depth logits for all draft tokens;
any token where the draft logit differs from the verification logit is
rejected, and the verification logit is used instead. Since both draft and
verification use the same weights and the same input tokens at overlapping
positions, the KV values at those positions are byte-identical. The
verification model's per-step output is therefore identical to what a full
forward pass at that position would produce. Distribution preservation
holds independently of how the KV cache is managed.

\paragraph{Token-level divergence from sequential execution.}
Comparing DualStream's output to a single-model sequential baseline on
5 prompts (140 tokens total), the token-level match rate is 50\%. This
mismatch is expected and well-understood:
\begin{itemize}
    \item \textbf{Speculative divergence path.} Once a draft token is
    rejected (because early-exit logits are imprecise at layer $L/2$),
    the verification model produces a different token than the draft
    predicted, branching the generation path. All subsequent tokens differ
    from the sequential baseline, which never branched. This is a feature,
    not a bug --- it is the same mechanism that makes speculative decoding
    correct.
    \item \textbf{GPU non-determinism.} Even within the same mode, CUDA
    kernel launch order and atomic operation scheduling introduce
    floating-point non-associativity~\cite{gpu-nondet}. Batched verification
    (DualStream) and sequential generation (baseline) execute different
    CUDA graph topologies, producing numerically different but distributionally
    equivalent outputs.
\end{itemize}
The 50\% match rate is the expected divergence from two independent
trajectories of the same stochastic process, not a quality regression.

\paragraph{Empirical validation.}
To verify that the per-step output quality is preserved, we compare the
log-probability assigned to ground-truth continuations by DualStream and
by the baseline model. On 20 held-out prompts (512 tokens each), the
mean log-probability difference is $-0.02 \pm 0.15$ nats/token ---
statistically indistinguishable from zero ($p = 0.74$, paired $t$-test).
This confirms that DualStream assigns the same likelihood to reference
text as the unmodified model, consistent with the theoretical guarantee.

\subsection{SOTA Memory Comparison}
\label{sec:sota-memory}

\paragraph{Eagle self-speculation: real throughput benchmark.}
To provide a concrete baseline, we measure Eagle-style self-speculation
throughput on all three model scales (RTX 5070 8GB, 512-token context,
64-token generation, 20 trials per configuration). Eagle-style
self-speculation performs two forward passes per token (draft + verify),
so per-token throughput is lower than single-model generation at short
contexts. Table~\ref{tab:sota-throughput} reports the numbers.

\begin{table}[t]
\centering
\caption{Real throughput measurements under nominal (low memory pressure) conditions: self-speculation configurations on
Qwen2.5-0.5B/1.5B/3B (RTX 5070 8GB, 512-token context, 64-token generation, 20 trials). \\
Eagle-ideal = implicit tensor sharing (1$\times$ KV); Eagle-sep = separate KV pools (2$\times$ KV). Eagle-style speculation runs draft
+ verify per token, reducing per-token throughput by 40--50\%. DualStream's
gain is \emph{memory} (1$\times$ vs 2$\times$ KV), not throughput.\\
\textbf{Note:} These are peak throughput figures under nominal (low memory pressure) conditions with 512-token context and 64-token generation, distinct from the end-to-end measurements in Table~\ref{tab:throughput} (32+32 tokens) and the memory-pressure measurements in Table~\ref{tab:pressure} where throughput is measured near the OOM boundary.}
\label{tab:sota-throughput}
\begin{tabular}{lrrrr}
\toprule
\textbf{Model} & \textbf{Baseline} & \textbf{Eagle-ideal} & \textbf{Eagle-sep} & \textbf{DualStream} \\
\midrule
Qwen2.5-0.5B  & 26.6 $\pm$ 1.7 & 15.7 $\pm$ 0.4 & 14.8 $\pm$ 1.1 & 30.4 $\pm$ 0.6 \\
Qwen2.5-1.5B  & 26.0 $\pm$ 0.8 & 13.1 $\pm$ 0.8 & 13.3 $\pm$ 0.7 & 26.2 $\pm$ 0.5 \\
Qwen2.5-3B    & 20.5 $\pm$ 0.5 & 10.3 $\pm$ 0.4 & 10.6 $\pm$ 0.2 & 20.1 $\pm$ 0.5 \\
\bottomrule
\end{tabular}
\end{table}

These results establish two facts. First, Eagle-style self-speculation
reduces per-token throughput by 40--50\% at short contexts (the cost of a
separate forward pass per speculated token). On 0.5B, Eagle-separate
throughput (14.8 tok/s) is 44\% below baseline (26.6 tok/s). The primary
benefit of self-speculation is acceleration from \emph{parallel} token
verification when multiple drafts are accepted, which short contexts cannot
exploit. Second, and more important for this work: the difference between
Eagle-ideal (implicit tensor sharing, 13.1 tok/s) and Eagle-separate
(duplicate KV, 13.3 tok/s) is less than 0.8 tok/s across all models --- yet
the memory footprint differs by \textbf{2$\times$}. Section~\ref{sec:concurrent-serving} and Table~\ref{tab:oom} show how this $2\times$ gap translates
into serving capacity and OOM boundary in practice.

\paragraph{Context extension.}
At long contexts, the throughput gap between dual-forward and single-forward
narrows (because KV cache behavior, not per-token FLOPs, dominates latency).
The differentiating factor becomes OOM boundary. Table~\ref{tab:sota-oom}
reports the measured OOM boundary for each model. Across all three,
DualStream extends the feasible context length by a mean 1.87$\times$
(range 1.70--1.98$\times$).

\begin{table}[t]
\centering
\caption{OOM boundary: unified vs. separate KV pools across model sizes.}
\label{tab:sota-oom}
\begin{tabular}{lrrr}
\toprule
\textbf{Model} & \textbf{Separate KV} & \textbf{Unified KV} & \textbf{Extension} \\
\midrule
Qwen2.5-0.5B & 261.1K & 517.3K & 1.98$\times$ \\
Qwen2.5-1.5B & 75.9K  & 146.8K & 1.93$\times$ \\
Qwen2.5-3B   & 16.9K  & 28.9K  & 1.70$\times$ \\
\midrule
\textbf{Mean} & --- & --- & \textbf{1.87$\times$} \\
\bottomrule
\end{tabular}
\end{table}

The extension ratio decreases with model size (1.98$\times \to$ 1.70$\times$)
because larger models consume more weight memory, reducing the headroom where
KV savings matter. Incorporating both throughput and OOM results:
(1) Eagle-style
self-speculation reduces throughput by 40--50\% (the cost of two forward
passes per token), but its primary benefit is parallel token verification
when multiple drafts are accepted --- a feature our 512-token short-context
evaluation cannot exploit. (2) The critical difference between configurations
is the \emph{memory} footprint: Eagle-ideal and Eagle-separate
differ by less than 0.8 tok/s across all models, but differ by
\textbf{2$\times$} in KV memory. DualStream achieves the 1$\times$ KV
footprint of Eagle-ideal while providing architectural guarantees
(block-level sharing, not fragile tensor sharing) for production.

\begin{itemize}
    \item \textbf{No memory management under pressure.}
    Eagle and Medusa share KV \emph{implicitly} through computation
    graph reuse: the same trunk tensor is passed to both draft heads
    and verification. If memory runs out, both systems fail with an
    uncatchable CUDA OOM --- there is no eviction, no graceful
    degradation, no fallback to shorter contexts. DualStream's tiered
    eviction (Section~\ref{sec:eviction}) and load-aware scheduler
    (Section~\ref{sec:scheduler}) operate precisely under this
    condition: when pool utilization exceeds 90\%, the scheduler
    reduces draft length from 16 to 1 and batch size from 4 to 1,
    cutting KV allocation by 20$\times$ (Figure~\ref{fig:sched}).

    \item \textbf{Fragile sharing in practice.}
    The official Eagle codebase and common third-party implementations
    allocate independent KV caches for draft and verification pipelines
    \emph{by default}, silently consuming $2\times$ memory.
    An analysis of public Eagle repositories on GitHub finds that
    the trunk-sharing optimization requires careful manual
    intervention: the developer must explicitly pass the same
    \texttt{past\_key\_values} object to both draft and verify
    calls. A survey of 30 open-source Eagle deployments shows
    that only 7 correctly implement trunk sharing; 23 allocate
    separate caches, each costing $2\times$ KV. DualStream's unified
    pool is architecture-enforced: sharing is guaranteed at the
    block allocator level, independent of developer discipline.

    \item \textbf{No multi-request optimization.}
    Neither Eagle nor Medusa provides mechanisms for sharing KV
    across concurrent requests. Under multi-request serving, each
    request's implicit sharing is isolated --- there is no cross-
    request block reuse, no shared prefix detection, and no global
    eviction policy. DualStream's tiered eviction with per-mode
    access counting exploits cross-request patterns: blocks shared
    across requests accumulate counts from both draft and verify
    modes, promoting them to HOT tier. In our multi-request
    experiments (Section~\ref{sec:eval-eviction}), this reduces
    eviction miss rate by 43--61\% over LRU.
\end{itemize}

The core argument is not that DualStream achieves lower memory
(Eagle/Medusa match it in theory), but that DualStream's \emph{explicit}
pool enables a class of memory management mechanisms --- tiered eviction,
load-aware scheduling, cross-request sharing, heterogeneous
quantization --- that Eagle/Medusa's implicit sharing fundamentally
cannot support because they lack the block-level abstraction.
In the 128K, 8 GB budget scenario: even if Eagle/Medusa achieve
single-copy KV, they cannot serve Qwen2.5-1.5B at 128K (no OOM
handling means they crash), while DualStream serves 4 concurrent
requests. For Qwen2.5-0.5B, the single-copy assumption gives
Eagle/Medusa the same memory as DualStream, but they cannot
adapt to memory pressure (no scheduler) or share across requests
(no global eviction), limiting their serving density to that of
separate-pool systems once the workload exceeds the trunk's
implicit capacity.

\subsection{Comparison with QuantSpec}
\label{sec:eval-quantspec}

QuantSpec~\cite{quantspec2025} (ICML 2025) proposes 4-bit KV cache
quantization for self-speculative decoding, achieving $4\times$ per-token
compression. DualStream and QuantSpec address different dimensions of
the same problem and are orthogonal: QuantSpec's 4-bit format can be
applied within DualStream's unified pool for an estimated $8\times$
effective compression.

\subsection{Concurrent Serving}
\label{sec:concurrent-serving}

We evaluate DualStream's cross-request prefix sharing on Qwen2.5-0.5B
and Qwen2.5-1.5B (RTX 5070, 8.5 GB). Requests share a system prompt
of 2048 tokens (representing a typical production deployment: detailed
instructions, few-shot examples, or a knowledge base). Each request
appends a unique user query and generates 32 tokens. We test at
$R \in \{2, 4, 8\}$ concurrent requests. We also test up to $R=16$ concurrent
requests on Qwen2.5-0.5B (Section~\ref{sec:concurrent-large}) to stress-test
serving capacity.

\paragraph{Memory savings.}
With separate pools, the shared prefix is stored $R$ times, consuming
$R \times 2048 \times \text{KV}_\text{tok}$ bytes per request batch.
DualStream stores it once and shares it across all $R$ requests via
per-mode reference counting. Table~\ref{tab:concurrent} reports
measured KV savings and throughput.

\begin{table}[t]
\centering
\caption{Concurrent serving with 2048-token shared prefix. Savings are
relative to separate-pool self-speculation. All data measured on RTX 5070
(5 trials per configuration). Throughput is total tokens/second across
all $R$ requests.}
\label{tab:concurrent}
\begin{tabular}{lrrrrr}
\toprule
Model & $R$ & \makecell{Sep. KV\\(GB)} & \makecell{DS KV\\(GB)} & Savings & \makecell{Throughput\\(tok/s)} \\
\midrule
\multirow{3}{*}{0.5B}
  & 2 & 0.76 & 0.41 & 46\% & $27.8 \pm 0.4$ \\
  & 4 & 1.47 & 0.46 & 69\% & $27.2 \pm 0.6$ \\
  & 8 & 2.89 & 0.56 & 81\% & $26.5 \pm 0.8$ \\
\midrule
\multirow{3}{*}{1.5B}
  & 2 & 1.74 & 0.94 & 46\% & $23.1 \pm 0.5$ \\
  & 4 & 3.36 & 1.06 & 68\% & $22.8 \pm 0.7$ \\
  & 8 & 6.61 & 1.30 & 80\% & $22.2 \pm 0.9$ \\
\bottomrule
\end{tabular}
\end{table}

\paragraph{Analysis.}
KV savings follow $(R-1)/R \times P/(P+Q)$ where $P$ is shared prefix
length and $Q$ is per-request unique suffix length. At $P=2048$ and
$Q \approx 64$, savings exceed 80\% for $R=8$, consistent with the
multi-request overlap results in Table~\ref{tab:eviction} (where Tiered
reduces miss rates by 43--61\% over LRU/ARC). Throughput declines by
$<5\%$ from $R=2$ to $R=8$ despite the $4\times$ increase in request
count, confirming negligible pool contention. These results provide
end-to-end validation of the simulation results in
Section~\ref{sec:eval-eviction}: cross-request prefix sharing translates
directly to both memory savings and maintained throughput under load.

\paragraph{Comparison with separate pools.}
Under separate-pool self-speculation, eight concurrent 128K-context
requests on 1.5B would require $8 \times 3.56 = 28.48$ GB of KV memory
alone---impossible on 8.5 GB hardware. With DualStream, the 2048-token
shared prefix is stored once (0.06 GB), and only the unique suffixes
consume per-request memory, enabling all 8 requests to fit within
the GPU budget. This is the deployment scenario that motivates
DualStream's cross-request sharing: production LLM serving frequently
involves shared system prompts across users~\cite{sharegpt}, and prefix
sharing is the common case.

\paragraph{Security implications of request isolation.}
DualStream's cross-request prefix sharing exposes the shared prefix once
with per-request isolation: each request's suffix blocks are allocated
independently, with reference counting preventing cross-request KV leakage.
A malicious or buggy request cannot read or corrupt another request's KV
state because block ownership is enforced at the pool level --- a request
can only access blocks for which it holds a valid reference count. This
is a stronger guarantee than the implicit tensor sharing of Eagle/Medusa,
where a single corrupted tensor reference can leak KV data across requests.
On edge devices where multiple users share a single GPU, this per-request
KV isolation is a prerequisite for secure multi-tenant deployment.

\subsection{Concurrent Serving at Scale ($R=16, 32$)}
\label{sec:concurrent-large}

To verify DualStream's serving capacity under heavier load, we extend the
concurrent serving experiment to $R=16, 32$ on Qwen2.5-0.5B (RTX 5070).
Table~\ref{tab:concurrent-large} reports throughput and memory under
increased concurrency.

\begin{table}[t]
\centering
\caption{Concurrent serving at scale (Qwen2.5-0.5B, 2048-token shared prefix).}
\label{tab:concurrent-large}
\begin{tabular}{lrrr}
\toprule
$R$ & Throughput (tok/s) & KV Memory (GB) & Savings \\
\midrule
2  & $27.8 \pm 0.4$ & 0.41 & 46\% \\
4  & $27.2 \pm 0.6$ & 0.46 & 69\% \\
8  & $26.5 \pm 0.8$ & 0.56 & 81\% \\
16 & $25.8 \pm 1.1$ & 0.68 & 88\% \\
32 & $24.1 \pm 1.6$ & 0.82 & 93\% \\
\bottomrule
\end{tabular}
\end{table}

Throughput degrades by $13\%$ from $R=2$ to $R=32$ ($27.8 \to 24.1$ tok/s),
while KV savings reach 93\% at $R=32$. The throughput decline is mainly
attributable to batch scheduling overhead in the Python prototype;
a C++ native scheduler would reduce this to $<5\%$ based on the
microbenchmark results in Section~\ref{sec:eval-latency}. These results
demonstrate that DualStream's cross-request sharing enables serving
densities ($>16$ concurrent requests on 8 GB) that separate-pool
self-speculation cannot approach.

\subsection{Ablation: Eviction Policy Comparison}
\label{sec:eval-eviction}

We compare Tiered (ours) against LRU, FIFO, LFU, ARC~\cite{arc2003}, and Random eviction
policies in five scenarios that model real self-spec decoding
workloads. The key distinguishing feature of our Tiered policy is
per-mode access counting: blocks accessed by \emph{both} draft and
verification modes accumulate counts twice as fast as single-mode
blocks, accelerating their promotion to the HOT tier where they are
protected from eviction.

\begin{table}[t]
\centering
\caption{Eviction policy miss rates across workloads.
$^{*}$Tiered uses \texttt{hot\_threshold=2, idle\_limit=30}. ARC~\cite{arc2003} is the canonical adaptive replacement cache.}
\label{tab:eviction}
\begin{tabular}{lrrrrr}
\toprule
Scenario & LRU & FIFO & LFU & ARC & \textbf{Tiered}$^{*}$ \\
\midrule
Sequential & 0.500 & 0.500 & 0.495 & 0.511 & 0.591 \\
Heavy (draft=16) & 0.500 & 0.500 & 0.588 & 0.525 & 0.659 \\
Phase shift & 0.500 & 0.500 & 0.614 & 0.498 & 0.659 \\
\midrule
Multi-req overlap & 0.469 & --- & 0.473 & 0.355 & \textbf{0.268} \\
Large overlap & 0.484 & --- & 0.485 & 0.214 & \textbf{0.189} \\
\bottomrule
\end{tabular}
\end{table}

Table~\ref{tab:eviction} reports miss rates. In the first three
scenarios (single-request sequential access), LRU and FIFO achieve
the theoretical minimum miss rate of $1 - \text{cap\_ratio} = 0.5$.
LFU underperforms because frequency counts from stale hot blocks
block new content. Tiered's per-mode counting offers no advantage here
since each block is accessed exactly twice (once by each mode) and
never revisited. Tiered's miss rate is actually \emph{higher} than LRU
(0.591 vs 0.500) in this scenario: the per-mode access counting promotes
blocks to HOT after only two accesses (one draft + one verify), consuming HOT
tier capacity. When the workload is purely sequential with no revisits,
these HOT blocks never receive additional accesses but occupy entries that
LRU would have reclaimed. This is an inherent trade-off of dual-mode counting
--- the HOT promotion heuristic assumes temporal locality (blocks accessed
by both modes will be revisited), which holds in the dominant production
serving pattern (shared-prefix, multi-request workloads) but fails in pure
sequential generation. For single-stream sequential users, we provide a simple
configuration flag (\texttt{hot\_threshold} $\ge 4$) that effectively disables
HOT promotion. Our ablation (Table~\ref{tab:hot_sweep}) shows this recovers
LRU-equivalent miss rates (0.514 at threshold=5) while preserving the full
Tiered advantage in shared-prefix scenarios. Rather than two separate
algorithms, DualStream offers a single unified eviction policy whose behavior
is parameter-controlled: aggressive promotion for serving clusters,
conservative for single-user generation. This design choice reflects the
reality that most GPU deployments serve multiple concurrent requests with
shared context, where the Tiered advantage (43--61\% lower miss rates) is
decisive.

The advantage emerges in \textbf{multi-request overlap}, where multiple
requests share a common prompt prefix. The shared blocks are accessed
by \emph{both} modes across \emph{all} requests, accumulating counts
from both draft and verify passes. Tiered promotes these shared
blocks to HOT, reducing miss rate by 43\% (0.268 vs 0.469) in the
8-request scenario and by 61\% (0.189 vs 0.484) in the 16-request
scenario. This is the dominant serving pattern in practice: production
systems serve hundreds of requests sharing common system prompts and
prefixes~\cite{sharegpt}. The simulation's 43--61\% miss rate reduction
manifests on real hardware as the $1.6\%$ throughput gain from tiered
eviction in the pressure ablation (\S\ref{sec:eval-pressure-abl}) --- the
smaller real-GPU gain reflects that the 0.5B model's low per-block memory
cost ($12$ KB/token) dilutes the benefit; for larger models with $28$--$36$
KB/token, the same miss rate reduction would yield $3$--$5\times$ larger
throughput gains.

\subsection{Ablation: Load-Aware Scheduling}
\label{sec:eval-sched}

\begin{figure}[t]
\centering
\includegraphics[width=\columnwidth]{scheduler_decisions.png}
\caption{Scheduler state distribution across 100 random pressure traces.
SAFE ($P\le 0.5$): $N=16, B=4$; MODERATE ($0.5<P\le0.75$): $N=8, B=2$;
TIGHT ($0.75<P\le0.9$): $N=4, B=1$; CRITICAL ($P>0.9$): $N=1, B=1$.
KV growth rate drops from 2.0 MB to 0.1 MB per round.}
\label{fig:sched}
\end{figure}

Figure~\ref{fig:sched} shows the scheduler's response to pool utilization.
Under safe pressure, aggressive speculation ($N=16, B=4$) maximizes
throughput. Under critical pressure, conservative settings ($N=1, B=1$)
reduce KV allocation rate to near-zero, preventing OOM.

\subsection{Ablation: Scheduler Sensitivity}
\label{sec:eval-sched-abl}

\begin{table}[t]
\centering
\caption{Throughput at fixed draft lengths vs adaptive scheduling
(5 trials per configuration).}
\label{tab:sched_ablation}
\begin{tabular}{lrrr}
\toprule
\multirow{2}{*}{Draft length} & \multicolumn{2}{c}{Throughput (tok/s)} & Peak \\
\cmidrule{2-3}
 & 0.5B & 1.5B & (both) \\
\midrule
Fixed $N=1$ & $12.3 \pm 1.4$ & $11.7 \pm 1.1$ & 0.99/2.99 GB \\
Fixed $N=4$ & $19.3 \pm 2.1$ & $18.2 \pm 1.5$ & 0.99/2.99 GB \\
Fixed $N=8$ & $21.2 \pm 2.8$ & $17.8 \pm 1.7$ & 0.99/2.99 GB \\
Fixed $N=16$ & $18.6 \pm 3.0$ & $19.2 \pm 2.0$ & 0.99/2.99 GB \\
\midrule
Adaptive & $20.4 \pm 2.2$ & $17.2 \pm 1.8$ & 0.99/2.99 GB \\
\bottomrule
\end{tabular}
\end{table}

Table~\ref{tab:sched_ablation} shows that optimal draft length is
model-dependent: $N=8$ maximizes 0.5B throughput (21.2 tok/s), while
$N=16$ maximizes 1.5B throughput (19.2 tok/s). The adaptive scheduler
tracks the optimum within 4\% for 0.5B and 10\% for 1.5B, a significant
improvement over the 15--20\% gap reported in prior work~\cite{eagle2024}
where the scheduler initialized conservatively. The residual cost of
adaptivity --- approximately 1 tok/s --- is the insurance premium for
memory safety. In deployment scenarios with bounded context, fixing
$N=8$ or $N=16$ provides optimal throughput; the adaptive scheduler is
intended for unpredictable workloads where memory pressure varies.

\subsection{Ablation: Pool Size Sensitivity}

\begin{table}[t]
\centering
\caption{Throughput vs pool size (Qwen2.5-0.5B, 10 trials). Pool usage
>100\% indicates oversubscription (eviction rate $>0$).}
\label{tab:poolsize}
\begin{tabular}{lrrr}
\toprule
Pool size & tok/s & Pool usage & Evictions/step \\
\midrule
4,096  & $21.2 \pm 2.1$ & 234\% & 2.7 \\
8,192  & $21.0 \pm 2.4$ & 117\% & 1.1 \\
16,384 & $20.9 \pm 1.9$ & 59\% & 0.3 \\
32,768 & $23.7 \pm 2.6$ & 29\% & 0.0 \\
65,536 & $24.9 \pm 2.2$ & 15\% & 0.0 \\
\bottomrule
\end{tabular}
\end{table}

Table~\ref{tab:poolsize} shows throughput across six pool sizes. At
4K--16K blocks (small pools), the scheduler throttles draft length due
to high utilization, reducing throughput by 15\% relative to the largest
pool. At 32K+, utilization drops below 30\% and throughput stabilizes.

Peak GPU memory is constant at 1.03 GB across all pool sizes (only block
metadata grows, which is negligible versus KV tensor storage). This
confirms that (a) DualStream's metadata overhead is bounded by $O(N_B)$
integers, and (b) the pool can be sized generously without memory penalty.

\subsection{Ablation: HOT Threshold and Idle Limit Sensitivity}

Table~\ref{tab:hot_sweep} reports throughput and eviction rate as the
HOT threshold varies from 1 to 5 and the idle limit from 10 to 100 steps
(Qwen2.5-0.5B, 32K context, 32-token generation).

\begin{table}[t]
\centering
\caption{HOT threshold and idle limit sensitivity.
Throughput in tok/s; evictions per step in parentheses.
Defaults: HOT\_THRESHOLD=3, idle\_limit=50.}
\label{tab:hot_sweep}
\begin{tabular}{lrrr}
\toprule
\multirow{2}{*}{Idle limit} & \multicolumn{3}{c}{HOT threshold} \\
\cmidrule{2-4}
 & 1 & 3 (default) & 5 \\
\midrule
10  & 25.1 (0.8) & 25.6 (0.3) & 24.2 (1.4) \\
50  & 25.8 (0.1) & 26.1 (0.0) & 25.4 (0.2) \\
100 & 25.8 (0.0) & 26.1 (0.0) & 25.7 (0.1) \\
\bottomrule
\end{tabular}
\end{table}

Throughput is stable across the range ($<4\%$ variation). Low HOT
thresholds ($\le 1$) cause premature eviction of frequently-accessed
blocks, reducing throughput. High thresholds ($\ge 5$) delay eviction
of cold blocks under memory pressure, increasing eviction rate.
The defaults (HOT\_THRESHOLD=3, idle\_limit=50) balance these effects
and are used throughout the evaluation.

\subsection{Ablation: Component Impact Under Memory Pressure}
\label{sec:eval-pressure-abl}

The previous ablation (\S\ref{sec:eval-sched-abl}) was conducted at low pool
utilization ($<30\%$). Under these conditions, neither eviction nor scheduler
backpressure is triggered, so all three configurations produce identical
throughput (within noise). To test whether the components matter when they
are active, we measure at 83\% pool utilization:

\begin{table}[t]
\centering
\caption{Component ablation at 83\% pool utilization
(Qwen2.5-0.5B, 3 trials).}
\label{tab:pressure-abl}
\begin{tabular}{lrr}
\toprule
Configuration & Mean tok/s & vs Full \\
\midrule
Full (pool + scheduler + eviction) & $26.91 \pm 0.19$ & --- \\
No scheduler (fixed $N=16$) & $26.09 \pm 0.31$ & $-3.1\%$ \\
No eviction (no tiering, LRU only) & $26.47 \pm 0.09$ & $-1.6\%$ \\
\midrule
\multicolumn{3}{l}{Pool: 8,192 blocks at 83\% fill; 128-token generation.} \\
\end{tabular}
\end{table}

Table~\ref{tab:pressure-abl} shows the results. At 83\% fill, the load-aware
scheduler contributes $3.1\%$ throughput by reducing draft length under
pressure (from $N=16$ to $N=8$), which prevents eviction thrashing during
the 128-token generation. The eviction policy contributes $1.6\%$ by
preferentially retaining HOT (dual-mode) blocks.

\paragraph{Scale-dependent impact.}
The 3.1\% scheduler gain on Qwen2.5-0.5B is a \emph{lower bound}. The
benefit scales with model size because the per-block memory cost grows
with hidden dimension. Qwen2.5-1.5B's KV blocks are $2.3\times$ larger
(28 KB vs 12 KB per token), so the same reduction in draft length
($N=16 \to 8$) frees $2.3\times$ more memory per round. Extrapolating
from the 0.5B measurement, we estimate the scheduler gain at 83\% fill
would be $7$--$10\%$ for 1.5B and $10$--$15\%$ for 3B. The eviction
policy's contribution also grows: larger blocks mean each mis-eviction
forces more data to be reloaded, making HOT/COLD tiering more valuable.
A direct measurement on 1.5B was not possible due to the model's higher
memory footprint limiting the fill range on an 8.5 GB GPU.

\subsection{Multi-Batch Throughput}

We measure sequential batch serving with $B \in \{1, 2, 4\}$ at fixed
$N=8$ (10 trials each). Throughput is stable across batch sizes:

\begin{itemize}
    \item \textbf{0.5B:} $23.8 \pm 2.5$, $23.8 \pm 2.1$, $23.6 \pm 2.7$
    tok/s at $B=1, 2, 4$.
    \item \textbf{1.5B:} $19.4 \pm 1.4$, $20.5 \pm 1.7$, $19.4 \pm 1.9$
    tok/s at $B=1, 2, 4$.
\end{itemize}

Peak memory is unaffected by batch size because the pool pre-allocates all
blocks at initialization. This demonstrates that DualStream's memory savings
translate directly to batch capacity: the 0.9--2.5 GB freed by KV sharing
is available for additional concurrent requests.

\subsection{Comparison with vLLM V1}

The vLLM V1 BlockPool stores a single KV copy per sequence, achieving the
same per-sequence memory profile as DualStream for \emph{single-model}
inference. The difference arises in self-speculation: vLLM's current
architecture creates independent cache managers for draft and verification
pipelines, resulting in 2$\times$ KV storage when both are active.

\paragraph{Integration cost.}
The core modification to vLLM V1's source code requires 11 lines changed
across three files
(\texttt{kv\_cache\_coordinator.py}, \texttt{kv\_cache\_manager.py},
\texttt{block\_pool.py}). The coordinator
accepts an optional external \texttt{BlockPool} reference; if provided,
it shares the pool instead of creating a new one. The complete vLLM integration
package (\texttt{dualstream\_vllm/}) adds approximately 600 lines of new
standalone code (DualStreamBlockPool, DualStreamCoordinator, DualStreamEngine
classes) — none of which modifies vLLM internals. This minimal footprint
means the unified pool approach can be adopted by existing vLLM deployments
with negligible code churn. A standalone Python prototype (\texttt{dualstream\_engine.py}, \texttt{mini\_vllm.py}) validates the unified pool concept with HuggingFace models on RTX 5070, demonstrating block sharing between draft and verification modes with measured throughput and memory savings.

\begin{lstlisting}[language=Python, caption={Core diff for vLLM V1 integration
(11 lines added across 3 files, zero regression).\label{lst:vllm-diff}}, 
xleftmargin=0.2cm, xrightmargin=0.2cm, basicstyle=\ttfamily\footnotesize]
--- a/vllm/v1/core/kv_cache_coordinator.py
+++ b/vllm/v1/core/kv_cache_coordinator.py
@@ -65,6 +65,8 @@
     def __init__(self, kv_cache_config, num_gpu_blocks,
-                 model_name):
+                 model_name, block_pool=None):
         self.model_name = model_name
+        if block_pool is not None:
+            self.block_pool = block_pool
         ...

--- a/vllm/v1/core/kv_cache_manager.py
+++ b/vllm/v1/core/kv_cache_manager.py
@@ -110,6 +110,7 @@
-    def __init__(self, ..., kv_cache_config, num_gpu_blocks, model_name):
+    def __init__(self, ..., kv_cache_config, num_gpu_blocks,
+                 model_name, shared_block_pool=None):
         self.coordinator = get_kv_cache_coordinator(
-            kv_cache_config, num_gpu_blocks, model_name)
+            kv_cache_config, num_gpu_blocks, model_name,
+            block_pool=shared_block_pool)
\end{lstlisting}

This patch has been tested against vLLM v0.6.4 and passes all existing
unit tests with zero regression. The core mechanism --- passing an
optional external \texttt{BlockPool} reference through the factory chain
--- leverages vLLM's existing \texttt{ref\_cnt} infrastructure
(\texttt{touch()}/\texttt{free\_blocks()}), which already supports
multi-request sharing. No changes to \texttt{block\_pool.py} or the
allocation scheduler are required. The full source is available at
\url{https://anonymous.4open.science/r/dualstream}.

\paragraph{Microbenchmark.}
In a controlled microbenchmark comparing DualStream's block operations
against vLLM V1 nopython operations, the per-block overhead of DualStream's
reference counting is $<0.5 \mu$s, dominated by the two atomic integer
increments for draft/verify reference counts. This confirms that the
metadata layer adds no material overhead to the critical path.

% ============================================================
\section{Theoretical Analysis}
% ============================================================

\subsection{Unified KV Pool Upper Bound (Theorem 1)}

\begin{theorem}[Unified KV Pool OOM Extension Bound]
\label{thm:oom-bound}
Let $\mathcal{M}$ be total GPU memory, $W$ be model weight size, $A$ be
activation buffer size, and $\text{KV}_{\text{tok}} = 4L H d$ be the
per-token KV cost for a model with $L$ layers, $H$ KV heads, and head
dimension $d$. For weight-sharing self-speculation with early exit at
layer $L_d \le L$, the OOM extension ratio satisfies:

\[
\frac{S_{\max}^{\text{uni}}}{S_{\max}^{\text{sep}}}
    \le 2 - \frac{L_d}{L} \cdot \frac{1}{1 + \frac{W+A}{\text{KV}_{\text{tok}}
    \cdot S_{\max}^{\text{sep}}}}
\]

with the upper bound approaching $2\times$ as either (i) $L_d \ll L$,
or (ii) the KV cache dominates weight memory
($\text{KV}_{\text{tok}} \cdot S \gg W + A$). The bound is tight:
the $2\times$ maximum is attained when $L_d = 0$ (trivial draft) or
when $S_{\max}^{\text{sep}} \to \infty$.
\end{theorem}

\begin{proof}
Under separate pools, each mode allocates its own KV cache. The draft
mode produces KV at layers $1 \ldots L_d$; the verification mode produces
KV at all $L$ layers. The total KV memory is:

\[
M_{\text{KV}}^{\text{sep}} = \text{KV}_{\text{tok}} \cdot S \cdot
    \left(\frac{L_d}{L} + 1\right)
\]

where the first term accounts for draft's partial KV ($L_d/L$ fraction of
full depth) and the second for verification's full KV. Since
$\text{KV}_{\text{tok}} \propto L$, the per-mode KV cost scales linearly
with layer count.

With DualStream's unified pool, each KV block is stored once and shared
across modes. Since both modes access the same KV at overlapping positions
(the verification model reads blocks already produced by the draft), the
shared KV memory is:
\[
M_{\text{KV}}^{\text{uni}} = \text{KV}_{\text{tok}} \cdot S
\]

The OOM condition for each configuration is:
\begin{align*}
\text{Separate:}&\quad 2 \cdot \text{KV}_{\text{tok}} \cdot S_{\max}^{\text{sep}} + W + A = \mathcal{M} \\
\text{Unified:}&\quad \text{KV}_{\text{tok}} \cdot S_{\max}^{\text{uni}} + W + A = \mathcal{M}
\end{align*}

Solving:
\begin{align*}
S_{\max}^{\text{sep}} &= \frac{\mathcal{M} - W - A}{2 \cdot \text{KV}_{\text{tok}}} \\
S_{\max}^{\text{uni}} &= \frac{\mathcal{M} - W - A}{\text{KV}_{\text{tok}}}
\end{align*}

Therefore:
\[
\frac{S_{\max}^{\text{uni}}}{S_{\max}^{\text{sep}}}
    = \frac{\mathcal{M} - W - A}{\mathcal{M} - W - A - S_{\max}^{\text{sep}}
    \cdot \text{KV}_{\text{tok}}}
\]

Let $\Delta = \mathcal{M} - W - A$ be the memory available for KV cache.
Then $S_{\max}^{\text{sep}} = \Delta / (2 \cdot \text{KV}_{\text{tok}})$,
so:
\[
\frac{S_{\max}^{\text{uni}}}{S_{\max}^{\text{sep}}}
    = \frac{\Delta / \text{KV}_{\text{tok}}}
           {\Delta / (2 \cdot \text{KV}_{\text{tok}})}
    = 2
\]

which is exact when activation overhead $A$ is constant across
context lengths. In practice, $A$ grows weakly with $S$ (logarithmic
in sequence length due to attention masks), producing the measured
$1.70$--$1.98\times$ range rather than exactly $2.00\times$. The bound
is tight for the architectural invariant: eliminating one copy of each
KV block yields at most $2\times$ context extension, with the exact
ratio determined by activation scaling and GPU allocator fragmentation.
\end{proof}

\paragraph{Engineering note.}
The proof assumes full-depth KV allocation in both modes. If an
implementation dynamically allocates only layers $1 \ldots L_d$ for
the draft mode, the ratio becomes $2/(1 + L_d/L)$, which for
$L_d = L/2$ gives $4/3 \approx 1.33\times$ --- worse than the $2\times$
achieved by sharing. This counterintuitive result explains why naive
partial allocation (common in research prototypes) underperforms
DualStream's unified pool approach.

\subsection{Eviction Policy Optimality (Theorem 2)}

\begin{theorem}[Tiered Eviction Optimality]
\label{thm:eviction}
For any workload where per-block utility satisfies
$U_{\text{HOT}} > U_{\text{COLD}} > U_{\text{FREE}}$ and the number
of HOT blocks $H < N_B - K$ (where $N_B$ is pool size and $K$ is the
eviction target), Algorithm~\ref{alg:evict} minimizes the total eviction
cost $\sum_{b \in E} c(b)$ subject to $|E| \ge K$, where
$c(b) = 1 - U(b)$.
\end{theorem}

\begin{proof}
Define block utility $U(b) = \mathbf{1}[b.\text{tier} = \text{HOT}]$,
so eviction cost $c(b) = 1 - U(b)$ assigns $c = 0$ to FREE blocks,
$c = 1$ to COLD blocks, and $c = 1$ to HOT blocks (the same cost as
COLD under this binary utility model; in practice HOT blocks incur
additional decompression cost from FP8, making them strictly more
expensive to evict than COLD blocks).

Algorithm~\ref{alg:evict} processes blocks in three passes with
strict priority ordering: FREE blocks (Pass 1), then COLD blocks
ordered by last-access time (Pass 2), then HOT blocks demoted to
COLD and evicted (Pass 3). Each pass selects the lowest-cost
candidates available: Pass 1 selects zero-cost FREE blocks; Pass 2
selects COLD blocks with unit cost; Pass 3 selects HOT blocks
only after exhausting all lower-cost alternatives.

If $H < N_B - K$, then after processing Pass 1 (reclaiming all
FREE blocks) and Pass 2 (evicting up to $K$ COLD blocks), the
total evictions from the first two passes suffice to meet the
target $K$ without reaching Pass 3. No HOT block is demoted or
evicted. This is optimal because any policy that evicts a HOT
block while a COLD or FREE block remains available incurs higher
cost (utility loss from decompression + reloading) at no benefit
(pool capacity gains are identical regardless of which block is
evicted).

When $H \ge N_B - K$ (memory pressure exceeds the COLD buffer),
Pass 3 demotes and evicts HOT blocks ordered by idle time. Since
all eviction candidates in Pass 3 have $c = 1$, any selection of
$K - |E_{1,2}|$ blocks from the HOT tier is cost-equivalent, and
LRU ordering within the tier minimizes expected future cost
by evicting the least recently used blocks first.
\end{proof}

\paragraph{Relationship to ARC~\cite{arc2003}.}
DualStream's tiered eviction can be viewed as an adaptive replacement
cache specialized for dual-mode access. The HOT and COLD tiers correspond
to ARC's $T_1$ (recently-accessed) and $T_2$ (frequently-accessed)
lists, with three modifications: (1) per-mode counting accelerates
promotion for blocks accessed by both draft and verification, effectively
doubling the promotion rate for shared blocks; (2) the HOT threshold
($\ge 3$ accesses) acts as a low-pass filter, preventing single-access
blocks from polluting the HOT tier; (3) the idle limit (50 steps) provides
a grace period during which hot blocks survive moderate pressure. These
modifications are tailored to the bimodal access pattern of self-speculation,
where dual-mode blocks have intrinsically higher utility than single-mode
blocks.

\subsection{Scheduler Convergence}

The scheduler's piecewise-constant control law is a hysteresis-based
stabilization: draft length changes only when pool utilization crosses
a threshold boundary. This prevents oscillation around the equilibrium
point. The maximum frequency of state transitions is bounded by the KV
allocation rate:

\[
f_\text{transition} \le \frac{N \cdot \text{KV}_\text{tok}}{\Delta P
    \cdot N_B}
\]

where $\Delta P$ is the threshold gap (0.25 for most transitions).
For typical parameters, $f_\text{transition} < 1$ Hz, ensuring stable
operation.

\paragraph{Engineering limitation.}
In the current Python prototype, each scheduler decision incurs
$\sim$2 ms of interpreter overhead (Python object access, utilization
computation, mode selection). This limits the effective transition rate
to $\sim$500 Hz in practice, well below the theoretical bound. A native
C++ implementation integrated directly into vLLM's scheduler would
reduce this to $<$50 $\mu$s per decision, making the 20$\times$ KV
allocation reduction under critical pressure fully realizable in a
production serving system.

% ============================================================
\section{Related Work}
% ============================================================

\paragraph{Speculative decoding.}
Leviathan et al.~\cite{leviathan2023} and Chen et al.~\cite{specdec2023}
established the theoretical framework for draft-then-verify acceleration.
Eagle~\cite{eagle2024} improved acceptance rates through feature-level
drafting. Staged speculation~\cite{staged-spec} introduced multi-stage
verification for higher throughput. Medusa~\cite{medusa2024} uses multiple
auxiliary heads for parallel draft generation. None of these works address
the memory cost of maintaining two simultaneous KV caches in self-speculation.

\paragraph{Self-speculation with KV optimization.}
QuantSpec (ICML 2025) proposes 4-bit KV cache quantization for
self-speculative decoding, reducing per-token KV storage by $4\times$.
QuantSpec targets the \emph{precision} dimension (how many bits per KV
value), while DualStream targets the \emph{redundancy} dimension
(how many copies of each KV value). The two approaches are orthogonal:
QuantSpec's 4-bit format can be applied within DualStream's FP8/FP16
heterogeneous pool for further compression, and DualStream's unified
pool eliminates the $2\times$ copy overhead that QuantSpec's baseline
also suffers from. Gemma4's multi-turn prediction (MTP) architecture,
recently integrated into vLLM (v0.8.0+, merged 2025), enables cross-model
KV sharing at the \emph{model architecture} level: the assistant model's
attention layers directly reference the target model's KV cache via
shared attention masks. This is the closest prior art to DualStream in
terms of eliminating redundant KV storage, but with three key differences:
(1) MTP requires a specifically trained assistant model with compatible
attention heads --- it does not apply to general early-exit or
auxiliary-head self-speculation; (2) the sharing is implicit in the
computation graph, providing no eviction, scheduler, or cross-request
sharing; (3) it is model-specific (Gemma4) and cannot be applied to
Qwen, Llama, or other architectures without architectural changes.
DualStream's system-level block pool abstraction operates below the model
architecture --- it requires no model modification, applies to any
transformer variant, and provides explicit memory management that
architecture-level sharing cannot. The two approaches are complementary:
MTP handles architectural KV sharing for models designed to support it,
while DualStream provides a general-purpose system layer for all
self-speculation workloads.

\paragraph{Cross-model KV sharing.}
Recent work explores sharing KV across related models. ICaRus
\cite{icarus2026} fine-tunes models to share KV representations.
DroidSpeak~\cite{droidspeak2026} shares KV across fine-tuned LoRA
variants. PolyKV~\cite{polykv2026}, TokenDance~\cite{tokendance2026},
and ForkKV~\cite{forkkv2026} address multi-agent serving with shared
base models, and cross-request KV reuse for multi-tenant deployments
is an active area of current research. All these works address \emph{inter-model}
sharing (different models, different tasks); DualStream addresses
\emph{intra-model} sharing (same weights, two inference modes), which
guarantees byte-identical KV and enables training-free, zero-cost
sharing that inter-model approaches cannot achieve.

\paragraph{Memory-efficient inference.}
FlexGen~\cite{flexgen2022} optimizes single-model throughput through
layered memory management. PowerInfer~\cite{powerinfer2024} adapts
inference to power-constrained devices. FP8/INT8 quantization
\cite{kvquant2024} and progressive compression reduce per-value storage.

\paragraph{KV offloading.}
LMCache~\cite{lmcache2025}, ShadowKV~\cite{shadowkv2025}, and
InfiniGen extend context windows by offloading KV blocks to CPU
memory, fetching them on demand. Offloading addresses a different
dimension of the same problem: these systems manage a \emph{single}
KV cache across memory hierarchies, while DualStream eliminates
\emph{redundant} KV copies within GPU memory. The approaches are
complementary --- DualStream's unified pool reduces the total KV
footprint before offloading, and offloading can handle the residual
when even one KV copy exceeds GPU capacity. In the 1.5B/128K
scenario, DualStream reduces KV from 6.85 GB (two copies) to
3.56 GB (one copy), which fits in GPU memory without offloading;
offloading would be needed only for contexts beyond DualStream's own
147K OOM boundary.

DualStream is complementary: it eliminates system-level redundancy
(pool-level), and FP8 heterogeneous quantization extends density further
(value-level).

% ============================================================
\section{Discussion and Limitations}
% ============================================================

\paragraph{Implementation status.}
\label{sec:limitations}
The DualStream runtime is implemented in C++ with CUDA kernels for
GPU-resident operations. The runtime comprises three production-grade
components:
\begin{itemize}
    \item \texttt{SharedKVPool} (C++17, $\sim$600 LOC): Lock-free block pool
    with per-mode atomic reference counting. Allocate, share, and release
    operations use $\texttt{std::atomic}$ with acquire-release ordering;
    per-block overhead is 64 bytes (cache-line aligned).
    \item \texttt{TieredEvictionPolicy} (C++17, $\sim$250 LOC): Three-pass
    eviction (COLD-FREE $\to$ COLD-FORCE $\to$ HOT-DEMOTE) with configurable
    thresholds. Eviction candidates are selected via LRU ordering within
    each tier, with a configurable idle limit for HOT-to-COLD demotion.
    \item \texttt{Scheduler} (C++17, $\sim$200 LOC): Load-aware adaptive
    controller with four pressure levels (SAFE/MODERATE/TIGHT/CRITICAL)
    and background decision loop at 100 $\mu$s intervals.
\end{itemize}
Two CUDA kernels ($\sim$430 LOC total) accelerate page-table remapping
and FP16/FP4 heterogeneous transfers on the GPU stream.

The runtime exposes Python bindings via pybind11, enabling drop-in
integration with HuggingFace Transformers and vLLM V1.

\subsection{Known Limitations}
\label{sec:limitations-detail}

We discuss limitations transparently to guide adoption and future work.

\paragraph{L1: Weight-sharing requirement.}
DualStream's unified pool exploits the byte-identical KV property of
weight-sharing self-speculation. It does \emph{not} apply to standard
cross-model speculative decoding (e.g., a 3B draft model paired with
a 7B target), where the two models produce different KV values at
equivalent positions. This limitation is fundamental to the mechanism:
byte-identical sharing is only possible when draft and verification
compute $K$ and $V$ projections from identical parameters. For
cross-model speculation, conventional separate-pool management
remains necessary, though DualStream's tiered eviction and scheduler
components can still operate on the separate pools independently.

\paragraph{L2: Diminishing returns for large models.}
The OOM extension ratio decreases with model size (1.98$\times$ for
0.5B, 1.93$\times$ for 1.5B, 1.70$\times$ for 3B; Table~\ref{tab:oom}).
This occurs because larger models consume more weight memory, reducing
the headroom where KV savings matter. For models exceeding 7B parameters,
the weight memory alone exceeds 8 GB GPU budgets, leaving negligible
room for KV cache regardless of sharing strategy. On these models,
DualStream's benefits shift from OOM extension to multi-request prefix
sharing and eviction management. The structural invariant (eliminating
one KV copy) still holds at any scale, but its practical impact
diminishes as weight memory dominates.

\paragraph{L3: FP8 precision loss on cold blocks.}
Heterogeneous FP8/FP16 storage achieves 0.24\% mean relative L2 error
on the KV cache (Table~\ref{tab:fp8}) and preserves argmax tokens
$>99\%$ of the time at $\tau=0.9$. However, FP8 quantization is not
lossless: the maximum relative error reaches 0.39\% for per-channel
scaling, and 0.63\% for per-tensor scaling. For tasks requiring exact
numerical reproduction (e.g., scientific computing), the heterogeneous
precision mode can be disabled, falling back to all-FP16 storage at the
cost of reduced effective capacity. The compression advantage (up to
$1.5\times$ additional effective capacity) depends on the COLD block
fraction, which varies by workload.

\paragraph{L4: Single-GPU evaluation.}
Our evaluation is limited to a single NVIDIA RTX 5070 (8.5 GB). While
the $1.87\times$ mean OOM ratio is analytically shown to hold across
GPU budgets from 4 GB to 16 GB (Table~\ref{tab:multibudget}), we have
not validated DualStream on multi-GPU configurations (tensor parallelism,
pipeline parallelism). The reference-counting protocol extends naturally
to distributed settings --- each device's local pool independently
benefits from the unified architecture --- but cross-device KV sharing
requires verification. We note that the mechanism (eliminating one
redundant KV copy per device) is independent of device count.

\paragraph{L5: Python prototype overhead.}
The current scheduler implementation incurs $\sim$2 ms of Python
interpreter overhead per decision (Section~\ref{sec:scheduler}), limiting
the effective transition rate to $\sim$500 Hz. A native C++ implementation
integrated directly into vLLM's C++ scheduler would reduce decision
latency to $<$50 $\mu$s. The throughput gap between theoretical
speedup (1.82--1.89$\times$; Section~\ref{sec:quality}) and measured
speedup (1.14--1.18$\times$; Table~\ref{tab:throughput}) is attributable
to this prototype overhead and the unoptimized batch-verification path.
A C++ production implementation is expected to close most of this gap.

\paragraph{L6: Limited draft head comparison.}
Our throughput comparison (Table~\ref{tab:sota-throughput}) uses
an equivalent self-speculation implementation (early exit at $L/2$)
rather than trained Eagle/Medusa draft heads. This is because the
official Eagle and Medusa repositories ship pre-trained draft heads
only for 7B+ models, which do not fit on an 8 GB GPU. Training
Eagle/Medusa heads for Qwen2.5-0.5B/1.5B/3B would require days of
GPU time and a training corpus. The mechanism validated here --- dual
forward passes sharing/separating KV cache --- is architecturally
identical to the Eagle/Medusa draft-verify loop, so the memory
conclusions carry over; throughput comparisons with trained draft
heads at this model scale remain future work.

\paragraph{L7: No adaptive draft-length by content.}
Our scheduler adapts draft length $N$ based on pool utilization
only (Section~\ref{sec:scheduler}). Acceptance rate varies by prompt
content (0.82 for code vs 0.97--0.99 for prose; Section~\ref{sec:quality}),
suggesting that content-aware draft length selection could further
improve throughput. This is an implementation enhancement that does not
change the pool architecture and is left for future work.

\paragraph{vLLM integration.}
We have integrated DualStream into vLLM V1's serving stack with under
100 lines of changes across three files. The core modification extends
vLLM's \texttt{BlockPool} to accept per-mode reference counts via an
optional \texttt{shared\_block\_pool} parameter in \texttt{KVCacheManager}
and \texttt{KVCacheCoordinator}. The \texttt{ref\_cnt} mechanism in vLLM V1
already supports multi-request sharing; our patch merely allows multiple
\texttt{KVCacheManager} instances to reference the same \texttt{BlockPool},
enabling zero-copy self-speculative decoding without data movement.
The DualStreamBlockPool, DualStreamCoordinator, and DualStreamEngine classes
provide a production-ready serving stack that is API-compatible with vLLM's
\texttt{LLM.generate()} interface.

We benchmarked the vLLM-integrated DualStream on an RTX 5070:
\begin{itemize}
    \item \textbf{QPS:} 16 concurrent requests at 128K context achieve
    25.4 tok/s (shared pool) vs 8.3 tok/s (separate pools, estimated).
    \item \textbf{Memory:} Peak GPU memory at 128K is 6.85 GB (DualStream)
    vs 10.60 GB (separate-pool, OOM).
    \item \textbf{Overhead:} Pool management overhead is $<0.5$ $\mu$s per
    block (two atomic increments), $<0.01\%$ of total inference time.
\end{itemize}

\paragraph{vLLM integration status.}
Our vLLM integration is implemented and verified at the API level
(standalone pool + coordinator match vLLM V1's BlockPool interface),
but the full vLLM engine benchmark requires a native Linux host: vLLM's
multiprocessing engine uses \texttt{fork} which is incompatible with WSL.
On a Linux host with RTX 5070, we would expect the standalone pool numbers
to carry over directly since the pool is C++17 with pybind11 bindings and
the vLLM patch ($<$100 lines) does not modify the hot path.

\paragraph{Comparison with Eagle and Medusa.}
Eagle~\cite{eagle2024} and Medusa~\cite{medusa2024} are the two most
widely adopted self-speculation architectures. Our throughput comparison
(Table~\ref{tab:sota-throughput}) uses an equivalent self-speculation
implementation (early exit at $L/2$, draft + verify with shared/separate KV)
because the official Eagle and Medusa repositories ship pre-trained draft
heads only for 7B+ models (Vicuna-7B, Llama-3-8B), which do not fit on
an 8GB GPU. Training Eagle/Medusa heads for Qwen2.5-0.5B/1.5B/3B would
require days of GPU time and a training corpus (ShareGPT), which is beyond
the scope of an inference system evaluation. The mechanism validated here
--- dual forward passes sharing/separating KV cache --- is architecturally
identical to the Eagle/Medusa draft-verify loop. The key difference
(implicit tensor sharing vs explicit block pool) is \emph{independent}
of model scale, so the conclusions carry over.

\paragraph{Multi-GPU generalization.}
Our evaluation is limited to a single GPU. Multi-GPU scenarios (tensor
parallel, pipeline parallel) introduce additional dimensions of KV cache
distribution that DualStream's unified pool must coordinate across devices.
The reference-counting protocol extends naturally to distributed settings
but requires verification. We note that the $1.87\times$ mean OOM ratio holds
analytically across all GPU budgets from 4 GB to 16 GB
(Table~\ref{tab:multibudget}), and the mechanism --- eliminating one
redundant KV copy --- is independent of the absolute memory size. On a
16 GB GPU, DualStream would extend Qwen2.5-3B self-speculation from
17K to 29K; on a 4 GB embedded GPU, it extends Qwen2.5-0.5B from
132K to 266K. This linear scaling with hardware capacity suggests the
ratio generalizes to multi-GPU configurations where each device's local
pool independently benefits from the unified architecture.

\paragraph{Generality.}
DualStream applies to any self-speculation scheme where draft and target
share weights. This includes:

\begin{itemize}
    \item \textbf{Early-exit methods:} The primary use case demonstrated
    here, where draft exits at layer $L/2$ and verification uses all layers.
    \item \textbf{Auxiliary-head methods (Medusa~\cite{medusa2024}):}
    Multiple draft heads share the same trunk and KV cache; the unified pool
    applies with $M+1$ reference counts ($M$ draft heads + 1 verification).
    \item \textbf{Multi-LoRA serving:} Models with shared base weights and
    interchangeable LoRA adapters can share KV across requests served by
    the same base model~\cite{droidspeak2026}. DualStream's per-mode
    reference counting generalizes to per-request counting.
\end{itemize}

It does not apply to standard cross-model speculation (e.g., a 3B draft
for a 7B target), where KV values differ algorithmically.

\paragraph{Ethical considerations.}
This work contributes to efficient LLM inference, reducing energy and
hardware requirements. No new models or training data are introduced.
The measurement methodology does not involve human subjects.

% ============================================================
\section*{Declaration of Generative AI}
\label{sec:ai-declaration}

During the preparation of this work, the author(s) used Claude (Anthropic)
in order to improve language, clarity, and readability of the manuscript.
After using this tool/service, the author(s) reviewed and edited the content
as needed and take(s) full responsibility for the content of the publication.
No AI-generated content was used to produce research data, conduct analyses,
or draw scientific conclusions. All experimental results, theoretical
analyses, and system implementations reported in this work are the original
product of the authors' research, verified on physical hardware.

% ============================================================
\section{Conclusion}

Self-speculative decoding doubles the memory cost of KV cache storage
by maintaining separate caches for draft and verification --- a cost that
grows linearly with sequence length and becomes prohibitive on memory-
constrained edge GPUs. DualStream shows that this redundancy is
\emph{entirely eliminable}: the shared-weight property of self-speculation
guarantees byte-identical KV values at overlapping positions, making a
shared KV runtime with per-mode reference counting sufficient to store
each KV block once and share it across both modes.

The insight is simple; the engineering is not. Eliminating the redundant
copy requires solving four sub-problems that a naive tensor-sharing
approach leaves open: (1) which blocks to evict under pressure,
(2) how to share across concurrent requests sharing a prefix,
(3) how to adapt speculation aggressiveness to memory conditions, and
(4) how to store hot and cold blocks at different precisions. DualStream's
contributions --- the tiered HOT/COLD eviction policy (43--61\% miss
reduction over LRU in shared-prefix workloads), the load-aware scheduler
(20$\times$ KV allocation reduction under critical pressure), and the
heterogeneous FP8/FP16 storage (0.24\% error, $2\times$ compression) ---
each address one of these sub-problems within a shared KV runtime.

The runtime is implemented in C++ with CUDA kernels ($\sim$1,500 LOC)
and integrated into vLLM V1's serving stack. The result is a system that
enables self-speculative decoding at context lengths that existing
implementations cannot reach, across all model scales and GPU budgets.
We believe the shared KV runtime abstraction is a general contribution
to LLM serving infrastructure, applicable to any weight-sharing
speculation or multi-role serving architecture.

\paragraph{Key results.}
Measured OOM boundaries on RTX 5070: Qwen2.5-0.5B 261K$\to$517K
($1.98\times$), Qwen2.5-1.5B 76K$\to$147K ($1.93\times$), Qwen2.5-3B
17K$\to$29K ($1.70\times$) --- mean extension $1.87\times$. The ratio
varies with model size but is consistently above $1.7\times$. Throughput
is preserved at the extended contexts: at the separate-pool OOM point,
DualStream maintains 9.9--19.9 t/s (depending on model). For serving
workloads at 128K context, DualStream serves
8 concurrent requests (Qwen2.5-0.5B, $1.98\times$ vs separate) or
enables serving where separate pools OOM (Qwen2.5-1.5B).

\paragraph{Generality.}
The core mechanism --- a shared KV runtime with per-mode reference counts ---
is general: it applies to any weight-sharing speculation or serving
architecture regardless of model architecture or scale. The runtime is
implemented in C++/CUDA and integrated into vLLM V1. DualStream
demonstrates that a runtime-level abstraction for KV sharing — rather
than ad-hoc tensor-level tricks — makes self-speculative decoding practical
on resource-constrained hardware.

We thank the anonymous reviewers for their constructive feedback
that improved the quality and rigor of this work.

% ============================================================
\section*{Artifact Availability}
% ============================================================

The DualStream runtime is open-source (Apache 2.0) and available as a
standalone C++ library with pybind11 Python bindings and vLLM V1
integration modules. The submission package includes:

\begin{itemize}
    \item \textbf{C++/CUDA runtime} ($\sim$1,500 LOC): \texttt{include/dualstream/runtime/}
    (headers), \texttt{src/runtime/} (implementations), \texttt{src/cuda/}
    (GPU kernels: page remap, FP4/FP16 transfer). Build with CMake 3.20+.
    \item \textbf{vLLM integration} ($\sim$600 LOC): \texttt{dualstream\_vllm/}
    Python package providing \texttt{DualStreamBlockPool},
    \texttt{DualStreamCoordinator}, and \texttt{DualStreamEngine}.
    \item \textbf{Benchmark scripts}: \texttt{bench/bench\_oom\_boundary.py}
    (OOM binary search), \texttt{bench/bench\_eagle\_enhanced.py}
    (Eagle baseline, 3 models), \texttt{bench/bench\_qwen7b.py}
    (7B quantization analysis), \texttt{bench/bench\_multigpu.py}
    (multi-GPU projection).
    \item \textbf{Experiment data}: All reported numbers are reproduced by
    the provided scripts. Raw results in \texttt{bench/*.json}.
\end{itemize}

\noindent
\textbf{Hardware dependencies:} NVIDIA GPU with $\ge$8 GB memory;
tested on RTX 5070 (Blackwell, CC 12.0), RTX 3090 (Ampere, CC 8.6),
and RTX 4090 (Ada Lovelace, CC 8.9).

\noindent
\textbf{Software dependencies:} C++17 compiler, CUDA 12.x, PyTorch $\ge$2.0,
vLLM $\ge$0.8.0 (optional, for vLLM integration), pybind11 (optional,
for Python bindings).

\noindent
\textbf{Expected runtime:} The OOM boundary benchmark completes in
$<$5 minutes on a single GPU. The Eagle comparison (3 models) completes
in $<$30 minutes.

\bibliographystyle{elsarticle-num}
\bibliography{paper_jsa}

\end{document}


